% ============================================================================
%  附录：审校批注汇编（由 % [Tutor Note: ...] 自动汇编，纯文本呈现）
% ============================================================================
\chapter{审校批注汇编}\label{app:notes}

\begin{quote}
\small
本附录把全书正文中以 \texttt{\% [Tutor Note: ...]} 形式就近标注的译者批注集中列出，
按章、节顺序排列，便于单独查阅。批注分四类：易错预警、作者省略的证明、原文校勘、
物理图像说明。批注只代表译者的核对意见，不改变正文。为保证本附录在任何 TeX 环境下
都能编译，此处以纯文本呈现公式片段；完整 LaTeX 形式见各章源码中的注释。
\end{quote}

\section{概述}
\subsection*{模形式/自守形式 / 粗略想法}
\begin{quote}\small
式 (1.5) 中的 $\omega$k/2 在 k 为奇数时并不是通常意义下的线丛（需要 spin 结构， 或按自守因子的半整数权来理解）。作者在这里把“模形式 = 典范丛幂次的全纯截面”一笔带过， 严格处理需要引入覆盖群 Mp2(R) 或 metaplectic 结构；这是后续讨论半整数权模形式时的常见陷阱。
\end{quote}

\subsection*{模形式/自守形式 / 结构}
\begin{quote}\small
式 (1.8) 只是“换基域”的隐喻性图示，并不是说 OF 等于 Z 作基变换。 正确的陈述是：OF 是 F 在 Q 上的整闭包；对 F/Q 有限，F 是 SpecZ 上的有限平坦覆盖 （在分歧素点处退化）。把它读成“base change”会误导后续对分歧的几何理解。
\end{quote}

\begin{quote}\small
式 (1.11) 中的记号 R$\times$ r1 表示 r1 个 R 的直积，C$\times$ r2 同理； 这里 $\times$ 是“笛卡尔积”而非群乘法。对 F=Q(d) 可直接验证：实二次域给出两个实嵌入， 虚二次域给出一个共轭对。后续计算 MG,F 的维数/体积时，这个分解是第一道易错关口。
\end{quote}

\subsection*{Langlands 纲领与 TFT / 概览}
\begin{quote}\small
式 (1.20) 中的 GC 指 Langlands 对偶群的 C-点，通常写作 G(C)。 这里 G 是 Q 上的约化群，G 是其复对偶群；而 Gal(F/F) 是绝对 Galois 群。 对 G=GLn，G=GLn(C)，这正是通常的 n 维 Galois 表示。对一般 G， “G-值表示”需要把 Galois 表示与对偶群的中心/连通分量问题一起考虑。
\end{quote}

\begin{quote}\small
式 (1.22) 的物理来源是 Kapustin–Witten 对 4d N=4 超对称 Yang–Mills 的 S-对偶作拓扑扭转；A-侧对应自守/磁侧，B-侧对应谱/电侧。它不是已证明的定理， 而是组织猜想的原则。这里作者用“等价”一词，读者应理解为“conjectural equivalence of TFTs”；严格陈述需要指明边界条件、算子与范畴的层级。
\end{quote}

\subsection*{Langlands 纲领与 TFT / 算术拓扑}
\begin{quote}\small
原书此表“圆盘”一行印作 SpecFt[[t]]（见 p.15），从上下文看应为 SpecFq[[t]]：q 是基域 Fq 的基数，t 是形式参数。这属于原稿笔误，译文按数学上正确的 Fq[[t]] 处理；若读者对照原 PDF 会发现此差异，特此说明。
\end{quote}

\begin{quote}\small
式 (1.27)–(1.29) 描述的是局部域绝对 Galois 群的驯分歧商： 惯性群 Iv 的驯部分同构于 $\prod$$\ell$'$\neq$ pZ$\ell$'(1)，而野惯性群是 pro-p 群； 商 Z$\ell$Z 中的 Z$\ell$ 来自 $\ell$-进循环（$\ell$$\neq$ p），Z 由 Frobenius 生成。 原文在 (1.29) 前有一句残句 “This tells us that there is For p=1, ...”， 其中 “there is” 之后的文字被截断/错位；“p=1” 与 “p=-1” 也不是素数， 而是 Baumslag–Solitar 群的退化参数。译文按 BS(1,1)=Z²、BS(1,-1)=Klein 瓶的基本群处理， 并保留原文的表述方式。
\end{quote}

\subsection*{Langlands 纲领与 TFT / $\mathcal{B}$-侧}
\begin{quote}\small
原文此处印作 “The B-wide is about studying...”，其中 “B-wide” 显然是 “B-side” 的笔误（p.19）。译文按 B-side 处理。
\end{quote}

\section{谱分解}
\subsection*{什么是谱？}
\begin{quote}\small
这里伴随的方向容易记反。O 是反变函子，故其“右伴随”应理解为在 Ringop 与 Spacesop 上的左伴随；若忘记取 op，就会把 Spec 误认为左伴随。正确读法是：Spec,R 表示“以 R 为函数环的自由空间”，而 (2.1) 表明这种自由性对一切测试空间成立。另注意 MapSpaces 必须理解为 plain hom-set（0-截断），否则伴随性只在同伦意义下成立。
\end{quote}

\subsection*{什么是谱？ / 紧支撑连续函数}
\begin{quote}\small
记号 Cv(X) 并非通用记号，通常写作 C0(X)（contract 支撑/消失于无穷）。若 X 非紧，C0(X) 不含单位元，此时应在单位化代数 C0(X)+ 上取极大理想，或改用“本原理想/特征”，否则 (2.4) 的 HomC*(A,C) 会漏掉零同态之外的信息。作者在正文中默认的应是“含单位的交换 C*-代数”与紧 Hausdorff 空间，非紧情形要靠单位化补救；这是引用 Gelfand--Naimark 时最常见的隐藏约定。
\end{quote}

\subsection*{什么是谱？ / 代数几何}
\begin{quote}\small
(2.8) 的交换图应读成：Spec: cRingop 与“整体截面”函子 $\Gamma$ 的伴随关系在仿射对象上的像；竖箭头表示“仿射概形 $\hookrightarrow$ 局部环空间”的满子范畴包含。作者用图把 Spec 画成从 cRing 指向仿射概形，实际方向是反变的，读图时须注意箭头两端是范畴而非函子。
\end{quote}

\subsection*{谱分解}
\begin{quote}\small
印刷版 (2.14) 末尾多出一个 w，从上下文看是排版残留（应与 (2.11) 中带下划线的 V 记号无关，也不进入任何公式）。照录于此；正确写法应为 V=$\lambda$$\in$ XV$\lambda$。另注意 (2.12) 的支撑本应定义为 Spec,R 中的闭子集而非 X，此处 X 是作者对“底空间”的泛称；在有限集情形二者一致，但一般 R 上应说 Supp(v)$\subset$Spec,R。
\end{quote}

\subsection*{谱定理 / 代数几何版本}
\begin{quote}\small
(2.21) 用并集而非直和，是因为不同 $\lambda$ 上的支撑集合可能相交于 0（即同时在两处支撑的向量只能为 0）；对有限生成模，实际只有有限多个 $\lambda$ 使 V$\lambda$$\neq$ 0，所以该并是有限并、无收敛问题。(2.24) 的下标写成 $\lambda$$\in$Spec 是省略了 Spec,k[x]=A1；$\lambda$ 表示“以 $\lambda$ 为中心的广义特征空间”，即 (2.25) 中的 k[x]/(x-$\lambda$)$\ell$i 部分，注意它与 (2.16) 的 V$\lambda$（真特征空间）不同：后者是前者的商。
\end{quote}

\subsection*{谱定理 / 可测版本}
\begin{quote}\small
(2.32) 左右两端用同一字母 $\pi$A 表示两个不同对象：左边带下划线的是层 $\pi$A，右边是投影算子值测度（对可测集取投影算子）。这种“同一符号、两种范畴”的写法是本节最容易混淆处。另注意 (2.31) 的积分只在 A 自伴时才是实轴上的谱积分；若 A 非自伴（或非正规），应改在 C 上取谱测度，且 (2.34) 的 Spec(A) 不再自动是实子集。
\end{quote}

\subsection*{谱定理 / 同伦版本}
\begin{quote}\small
印刷版 (2.39) 右端写作 C*(X)，但若如此则右边恒等于 M，与开集 U 无关，不可能定义出层。正确形式应是 M(U)=M$\otimes$C*(X)C*(U)：用 U 上的上链把常数模 M “局部化”。这里照录原文并加校勘；读者可按修正后的公式理解后面的陈述。(2.35) 的嵌入也不是逐点单射，而是在拟同构意义下把 R-模视为 X 上的局部常值复形。
\end{quote}

\subsection*{谱定理 / 物理诠释：可观测量与态}
\begin{quote}\small
印刷版 (2.43) 右端明显有误：|$\psi$i$\rangle$ 是未配对的右矢，且与左端（正数）量纲不符；按 Jordan 谱分解应写成 (1)/($\langle$|$\rangle$)$\Sigma$$\lambda$,$\psi$i$\lambda$,$\langle$$\psi$i|$\rangle$$\langle$|$\psi$i$\rangle$，即 $\lambda$|$\langle$$\psi$i|$\rangle$|2。这里照录原文并标出。另注意 (2.42) 的期望值用 $\langle$|$\rangle$ 归一化，与通常“态已归一化”的约定不同；作者在全书中保持这一不归一化的约定，在后续出现 $\rangle$ 的地方须与此一致。
\end{quote}

\section{Fourier 理论与阿贝尔对偶}
\subsection*{特征标}
\begin{quote}\small
印刷版 (3.2) 写作 Aut(X)，但按上下文（g$\mapsto$ Tg 是 V 上的算子、且右侧为 End(V)）应为 Aut(V)。这是原文笔误，照录并标出。另注意 Aut(V)$\subset$End(V) 只是幺半群包含：G 的群同态经此包含成为到乘法幺半群的映射，这正是下一节 (3.6)--(3.8) 要用“幺半群而非群”作伴随的原因。
\end{quote}

\subsection*{有限 Fourier 变换}
\begin{quote}\small
当 G 非交换时 C G 非交换，Spec 无从谈起；因此本章最终必须限制到 Abel 群（作者随后即如此假设）。(3.10) 把 C G 写成“函数之和”，实际只在 G 有限时成立；一般 G 应写成有限支撑函数（即测度）的代数，这就是正文紧接着强调“应把 f 想成测度”的原因。
\end{quote}

\begin{quote}\small
(3.21) 的对偶性要成立必须限制在 Abel 群 G 上；Spec(C G,$\ast$) 中的 $\ast$ 提醒读者这里取的是卷积乘法。另注意 (3.22) 的 g 是指 g 在对偶群中对应的特征标 $\chi$$\mapsto$$\chi$(g)，而不是某个新变量；下一小节 (3.26) 的双对偶映射会把它讲清楚。
\end{quote}

\subsection*{有限 Fourier 变换 / 隐藏的对称性}
\begin{quote}\small
(3.31) 的复共轭不能丢：在有限 Abel 群上 $\chi$(g,t)=$\chi$(g,t)-1，它正是逆变换中“正交关系”的来源，也是 (3.32) 不含共轭的原因。若把 (3.31)(3.32) 视为互逆变换，还必须补上归一化因子 1/|G|（或把 f 视为测度、 f 视为函数），作者在这里用的是“测度--函数”约定，故两式表面上不对称。
\end{quote}

\subsection*{Pontrjagin 对偶}
\begin{quote}\small
对一般 LCA 群，G 要取紧开拓扑才成为拓扑群（Pontrjagin 定理）；(3.34) 只写了集合。“G 紧 $\Leftrightarrow$ G 离散”这一条（见 p. 35）正是紧开拓扑的直接推论，后文 (3.50)--(3.51) 用它解释 e2$\pi$ inx$\in$ L2 与 $\delta$n$\in$$\ell$2 的差别。
\end{quote}

\subsection*{Pontrjagin 对偶 / 群代数}
\begin{quote}\small
印刷版 (3.35) 的被积函数用了自由变量 h，而积分变量是 g，这在原文中即不一致。按 (3.18) 的推前定义应写成 $\int$g$\in$ Gf1(g)f2(g-1h),dg，或把积分变量改为 h。叉积次序（gh-1 还是 g-1h）取决于群乘法约定，翻译时照录原文，读者按 f1$\ast$ f2(h)=$\int$ f1(g)f2(g-1h)dg 校正即可。
\end{quote}

\begin{quote}\small
严格说 L1(G) 只是 Banach *-代数而非 C*-代数（C*-包络才得到群 C*-代数 C*(G)）；作者说“C*-代数世界”是口语化的说法。用 L1(G) 的好处是（对 Abel G）Gelfand 理论直接给出 G，无需先取包络。
\end{quote}

\begin{quote}\small
(3.43) 右端照印本写作 L2(G)，按 Fourier 变换的定义应为 L2(G)（Plancherel 定理：Fourier 变换是 L2(G)$\to$ L2(G) 的等距同构）。这是原文的记号笔误，特此标出。
\end{quote}

\subsection*{Fourier 变换}
\begin{quote}\small
(3.53) 与 (3.54) 的左边都印作 T，但二者说的是不同对象的对偶：环面的对偶是格 $\Lambda$（离散），格的对偶是环面 $\Lambda$$\otimes$ZU(1)。作者用同一个符号 T 是记号滥用；后文 (3.81)(3.124) 中的 T 一律指“对偶环面”。这是阅读本节最容易出错处。
\end{quote}

\begin{quote}\small
(3.63) 印刷版右端为 “SpecRR[d/dx=t]”，其中连续出现两个 R 是排版重复，应读作 SpecRR[d/dx=t]：即在实系数多项式环 R[d/dx] 上取 Spec，得到 Rt。这里 g=R（G=R 的交换 Lie 代数），g*=Rt 是“动量/频率”坐标所在的直线，与 (3.57) 中 t 的角色一致。
\end{quote}

\subsection*{在量子力学中}
\begin{quote}\small
(3.66) 取 pj=i,d/dqj（数学约定，$\hbar$ 已吸收进 p），此时 [pj,qk]=i$\delta$jk；而 (3.67) 又写成 $\xi$ f=f$\xi$+$\hbar$ f'，两式的归一化并不一致——若按 (3.66)，(3.67) 右端应为 i f'（或把 p 改回物理约定 -i$\hbar$$\partial$）。此外 (3.67) 中 f' 应理解为 $\xi$(f)，即沿 $\xi$ 的方向导数，而非对某个坐标的偏导。这是本节最容易混用的符号。
\end{quote}

\subsection*{Cartier 对偶}
\begin{quote}\small
点函子的 R 取遍所有交换环（不只是域），这是“概形论”代数群与“经典”代数群的关键差别：例如 Z 作为代数群只在此意义下才是群概形。(3.72) 中的虚线提升不是自动的，它等价于给定 G(R) 的群结构对 R 函子化；作者略去的“群运算与结构态射相容”条件通常还要求乘法、单位、逆都是簇态射。
\end{quote}

\begin{quote}\small
(3.85) 的并集写成 nSpec,k[t]/(tn) 是（形式）归纳极限而非普通概形：每个 Spec,k[t]/(tn) 是加厚点，其并集是形式概形 Ga，作为函子它把 R 送到 R 的幂零元集。作者在下一段用“t 幂零/接近 0”来直观解释这一点。另注意 (3.87) 单独作为一条编号公式只是重述 Ga=Spec,k[x]，真正的“特征标公式”是 (3.88)。
\end{quote}

\subsection*{Cartier 对偶 / Fourier 级数例子}
\begin{quote}\small
(3.97) 的右边写成 QC(Z) 看起来奇怪，但它正是 (3.94) 的对称版本：Rep(G)$\simeq$QC(G)，取 G=Gm 时 G=Z，故 Rep(Gm)$\simeq$QC(Z)。注意区分：QC(Z) 是 Z（作为离散概形）上的拟凝聚层，即 Z-分次向量空间；而 QC(Gm) 是 k[z,z-1]-模。两者在 Fourier 对偶下互换，不要与 (3.94) 的 QC(Gm) 混淆。
\end{quote}

\begin{quote}\small
(3.110) 的 C[Z] 是群代数（有限支撑函数），其 Gelfand/代数对偶是 O(C$\times$)=C[z,z-1]；这正是 (3.94) 的另一种写法。(3.111)--(3.112) 说明代数对象（多项式/Laurent 多项式）作为 L2 的稠密子空间嵌入解析对象，但反之不成立：解析完成包含了代数版本没有的极限点。这一“代数 $\subset$ 解析”的包含关系在第 3.7 节的 A-侧/B-侧对比中会以“有限支撑 vs 局部常值”的形式再次出现。
\end{quote}

\begin{quote}\small
题外话 1 是全章的方法论总结：同一“对偶群”集合 G 上可装三种不同结构（测度 / 拓扑 / 代数），取决于群代数取哪一类完备化。这不是三种不同的对偶，而是同一对象的三种“几何化身”；第 3.7 节的 A-侧（离散/代数）与 B-侧（紧/解析）之分正是这一点的体现。
\end{quote}

\subsection*{Pontrjagin--Poincaré 对偶}
\begin{quote}\small
(3.113) 的成立需要 M 紧可定向、G 取合适的（有限或紧）系数，并且对偶取 Pontrjagin 对偶（连续特征标）；(3.114) 的第一支用 G$\otimes$ G(1) 的求值配对，第二支用 Hn(M,U(1))(1)（紧可定向 n 维流形上的积分），最后一支即积分。若 M 非紧或系数带拓扑，需要改用 Borel--Moore 同调/紧支撑对偶，这是后续章节（电磁对偶的边界条件）会再遇到的细节。
\end{quote}

\subsection*{Pontrjagin--Poincaré 对偶 / 二维}
\begin{quote}\small
(3.118) 用的是 $\pi$1(X)$\to$ G 的表示/同态集合，但严格说 LocG X 是群胚（每个对象带自同构），(3.119) 才对同构类取 $\pi$0；|$\cdot$| 就是取 $\pi$0。(3.117) 中 S1 的底层空间是栈：LocC$\times$S1$\simeq$ C$\times$//Z（由绕数作用），其“点”比同构类多。
\end{quote}

\begin{quote}\small
原文 3.7.2 与 3.7.3 的标题相同（都是 “Two dimensions”），并非重复印刷，而是两个互补的陈述：(3.115) 取 (H0,H1)，(3.120) 取 (H1,H0)，即交换 A/B 两侧的同调次数。表 3.1 中两行 “2” 分别对应这两条。T-对偶的半径 R$\leftrightarrow$1/R 正是这里 H1(S1,Z) 的“格”与 H0(S1,C$\times$) 的“对偶环面”之间的互换在半径上的体现。
\end{quote}

\subsection*{Pontrjagin--Poincaré 对偶 / 三维}
\begin{quote}\small
题外话 3 末尾的 “GU(1)-丛” 是排版错误，应为 “U(1)-丛”：从线丛 L 出发，Hom(L,X$\times$C) 的纤维是 Hom(C,C)$\cong$C，但可逆元给出 C$\times$=U(1)C，故得到主 C$\times$-丛（即 U(1)-丛）。另注意原文把 “principal bundle” 与 “gauge theory” 混用，是课程口语传统，指的都是带联络的主丛。
\end{quote}

\subsection*{Pontrjagin--Poincaré 对偶 / 四维}
\begin{quote}\small
本节标题是“四维”，但 (3.132)--(3.138) 中出现的流形都写作 M3：按 3.7 开头声明“几何对象比场论维数低一维”，四维理论的 A-侧应是三维空间 M3 上的线丛/规范场（而 B-侧是 M 上的平坦 C$\times$-丛）。这里的 M 与 M3 指同一个流形，只是作者在 (3.133) 中省略了上标；[M,K(Z,2)] 的 M 也应读作 M3。这是原文记号不一致处，不影响结论。
\end{quote}

\section{电磁对偶}
\subsection*{经典场论 / 二维}
\begin{quote}\small
此处 $\star$ 的归一化约定很关键：$\star$$\star$=$\pm$1（黎曼 2-流形上 对 1-形式 $\star$2=+1，对 3-维 Lorentz 流形上的 2-形式 $\star$2=-1）。 作者在 (4.9) 中取 B=-$\star$3E、在 (4.10) 中取 E=$\star$3B， 正是为保证 ($\star$ F) 的分量恰为 (B,E)；若采用相反的 Hodge 星 约定，这两个符号会同时反号，R$\leftrightarrow$ 1/R 的结论不变。
\end{quote}

\subsection*{经典场论 / 三维}
\begin{quote}\small
原文这一段的定义域仍写作 $\Sigma$，但本小节标题与上下文都是三维， 因此这里应理解为 3-流形 M3；照原样保留记号 $\Sigma$，读者按 M3 理解即可。
\end{quote}

\subsection*{四维／Maxwell 理论入门}
\begin{quote}\small
(4.13)、(4.14) 右端的 “\#” 是作者留的占位符，原书从未给出它 的具体数值。按 Gauss 单位制，b$\Sigma$=$\int$$\Sigma$F=4$\pi$ QB,$\Sigma$， 即在 SI 有理化单位制之外出现 4$\pi$ 因子；这说明这里的 \# 实际是与单位制 约定有关的常数（例如 4$\pi$），而不是数学上的未知量。
\end{quote}

\begin{quote}\small
(4.13) 左边下标误作 $\sigma$（应随右端作 $\Sigma$），而 (4.19) 中积分域写作 $\sigma$、右端又保留未展开的 [Sigma]，这三处是同一处 排版疏漏；译文按原样保留 $\sigma$ 并在 (4.19) 处以 [$\Sigma$] 呈现。
\end{quote}

\begin{quote}\small
(4.20) 原文第二行印作 “[$\Sigma$](L)$\in$ H2(M,Z)”，多出的 “(L)” 显然是笔误：[$\Sigma$] 是曲面 $\Sigma$ 的基本类，与线丛 L 无关。 译文已按数学含义排为 [$\Sigma$]，并在批注中保留原文痕迹。
\end{quote}

\begin{quote}\small
(4.19) 中 (1)/(2$\pi$ i)$\int$ F 取值为整数，正是第一 Chern 类的 de Rham 代表；这里 F 的归一化（U(1) 的生成元取 2$\pi$ iZ） 决定了 Dirac 量子化条件，与 (4.9) 的 Hodge 星符号约定相互独立。
\end{quote}

\subsection*{量子场论}
\begin{quote}\small
原文这里把空间维数在 Md-1 与 M3 之间来回切换：若 时空是 Md-1$\times$R，则时间切片应是 Md-1 而非 M3；下一段的 (4.23) 又回到 M3（此时 d=4）。这是讲义随写随改留下的记号不一致， 不影响结论，阅读时统一取 d=4 即可。
\end{quote}

\begin{quote}\small
(4.27) 的两个短正合列本来在书中以交换图排版（竖直箭头连接 $\Omega$1/$\Omega$1Z 与 Cflat、以及 H2 与 $\Omega$2Z）；这里用 array 重排，箭头的连接关系与原图一致。 关键点是右端：H2(M,Z) 与 $\Omega$2Z(M) 之间的竖直箭头表示 “整系数形式的上同调类”这一比较映射，而非 Chern 类映射。
\end{quote}

\subsection*{量子场论 / Lagrangian 形式}
\begin{quote}\small
“无序算子”与下一段的“有序算子”相对：有序（local）算子作用在 场空间上的泛函上，而无序（disorder）算子改变所考虑的场空间本身（引入奇点）。 这一区分在第 4.5.1 节会以 't Hooft 算子（改变 c1，即无序型）与 Wilson 算子（依赖联络，即有序型）的形式再次出现，务必对照阅读。
\end{quote}

\subsection*{量子场论 / 时间演化}
\begin{quote}\small
例 26 中若 N 是 (n-1)-流形，则 N$\times$ I 应是 n 维配边； 原文写作 d 维，是 n 与 d 记号的混用（本段讨论一般的 d 维理论，但 例子取的是 d=n）。
\end{quote}

\begin{quote}\small
(4.35) 原书是一个交换图：F(M) 在上，两条斜箭头分别指向 F($\partial$ Min) 与 F($\partial$ Mout)。这里用 array 重排，保留原编号；箭头方向与原图一致（都是“限制到边界”）。
\end{quote}

\subsection*{量子 Maxwell 理论}
\begin{quote}\small
(4.40) 的归一化是 Maxwell 理论中最容易搞混的地方。通常写 S=14e2$\int$ F$\star$ F+i$\theta$8$\pi$2$\int$ F F； 作者把 Yang--Mills 系数记成 (g)/(2$\pi$ i)、把拓扑项系数记成 $\theta$， 这只是把 g,$\theta$ 吸收进定义。要点是：$\int$ F F 只依赖线丛的同构类， 因而在连续形变下不变（这就是它被称为“拓扑项”的原因）；在 U(1) 情形它等于 c12=p1 的积分，取整数值，故 $\theta$ 以 2$\pi$ 为周期。
\end{quote}

\begin{quote}\small
(4.43) 中出现的 “BH1(M,Z)” 在书中没有定义，B 记号很 可能是笔误或作者自己的速记。数学上应读作 T$\simeq$Hom(H1(M),R/Z) $\simeq$ H1(M,Z)$\otimes$ZR/Z；若坚持 U(1) 取离散拓扑，则 T 成为 一个无穷离散群，“环面”的说法只在通常拓扑下成立。
\end{quote}

\subsection*{量子 Maxwell 理论 / 算子}
\begin{quote}\small
(4.47) 的交换图在原文中由“顶空间 + 两个投影 $\pi$in, $\pi$out + 底部的两个 H”构成，底部水平箭头标出的是积分核 $\pi$out*(e-S/$\hbar$D,$\pi$*in(-))； 这里用 array 重排，保留了全部箭头与编号。“挖去 $\gamma$ 后c1=1” 意味着联络在环绕 S2$\times$ S1 上的积分非零，故不能延拓——这正是 't Hooft 算子作为“无序算子”的几何含义。
\end{quote}

\subsection*{量子 Maxwell 理论 / 电磁对偶}
\begin{quote}\small
(4.53) 这个“等式”只在忽略向量空间因子 V 的意义下成立：两边 的 L2 都含有同一个 L2(V) 因子，作者把它略去。真正的内容是 Pontrjagin 对偶 L2(T)$\simeq$ L2($\Lambda$)，以及 L2($\Lambda$)$\simeq$ L2(T)，即 Fourier 变换交换 $\Lambda$ 与 T。 把 $\Lambda$B$\times$ TE 与 $\Lambda$E$\times$ TB 对照 (4.54)，就能看出 “F 侧对角化 Wilson、$\star$ F 侧对角化 't Hooft” 的含义。
\end{quote}

\subsection*{二维中的阿贝尔对偶}
\begin{quote}\small
(4.55)--(4.56) 中“S1R”表示半径 R 的圆。T-对偶的 标准陈述是：半径 R 的圆上的自由玻色子等价于半径 1/R（在 $\alpha$'=1 的单位下）的圆上的自由玻色子，同时把动量模与绕数模互换。这里的分级 H1(S1,Z)=Z（绕数）与 H0(S1,U(1))（取值）正是这一互换的 代数表现。原文用双重尖括号 $\langle$$\langle$$\cdots$$\rangle$$\rangle$ 排版，是一个 排版习惯，不影响含义。
\end{quote}

\subsection*{拓扑量子力学／量子场论}
\begin{quote}\small
这里“杀死时间”的准确含义是：在 Q-上同调（即把 Q-恰当形式 视为零）中，H=[Q,Q*] 作用为零，于是时间演化 eiTH/$\hbar$ 在同调上 变成恒等。注意 [Q,Q*]=H 的符号依赖于 Q* 取 d* 还是 -d*； 若取 Q*=-d*，则右边多一个负号，但“在 Q-上同调上 H=0”的结论不变。
\end{quote}

\subsection*{拓扑量子力学}
\begin{quote}\small
原文末句作 “N=8 had 16”，时态 “had” 是笔误， 应为 “has”。另外“奇算子个数 = 2N”这一计数只在 N$\le$ 4 时与通常的物理约定一致；对 N=8，超荷通常 计入其反粒子共轭，实际独立超荷数会因维数约化约定而不同，阅读时以原文的 计数为准即可。
\end{quote}

\begin{quote}\small
(4.66) 中 R$\Gamma$(OX) 的写法只在 X 紧合（例如射影） 时才与 H0,*(X) 同构；一般复流形上应写作 H0,*$\partial$(X) 与 R$\Gamma$(OX) 的导出范畴对象，而不是相等。作者在此处把 “H0(H)” 与 Dolbeault 上同调等同，是按紧 Kähler 情形书写的。
\end{quote}

\subsection*{拓扑量子力学 / 拓扑 Maxwell 理论}
\begin{quote}\small
(4.69) 的下标 U(1) 指对联络的自同构群（每点的 U(1) 旋转， 即规范变换）取等变上同调。若只取通常上同调，H•(C(M)) 会因 每个连通分支上的 BU(1) 因子而多出 C[u]；等变化正是为了把规范自由度 正确地商掉，这一点在 (4.73) 中体现为显式的 BU(1) 因子。
\end{quote}

\begin{quote}\small
(4.72) 原文直接写作 C[H2(M,Z)]，跳过了 C[$\pi$0(C(M))] 中“$\pi$0 由 c1 给出”这一步：由于 C(M)$\to$$\Lambda$ 的纤维是连通的环面，$\pi$0(C(M))$\cong$$\Lambda$。 另外注意 C[$\Lambda$] 是群代数，与 (4.75) 中 C[TC] 的对偶关系正是 Fourier 变换。
\end{quote}

\begin{quote}\small
表 4.1 原文以图像排版，这里按可用内容重排为 tabular， 以便检索与缩放；表中第 3 行“'t Hooft 算子 vs Wilson 算子”两栏内各有两行 说明，重排时用 quad 手动换行，语义与原文一致。
\end{quote}

\subsection*{拓扑量子力学 / 缺陷}
\begin{quote}\small
(4.81)、(4.83) 中 C[Loc($\cdots$)] 表示“该空间上的函数空间” （或函数空间的导出版本），并不是群代数；只有当 Loc 是离散群时才与群代数 重合。作者在 (4.72) 中 C[H2(M,Z)] 用的是群代数含义，两处记号相同而 含义不同，这是本书记号的一个陷阱。
\end{quote}

\begin{quote}\small
本节两类缺陷的物理含义不同：类时 't Hooft 线是“沿一条时间方向的 线插入磁流”，它改变 $\pi$0（即 c1），对应 A-侧增加 Z；类时 Wilson 线是“沿一条时间方向的线做和乐测量”，它改变纤维的平凡化，对应 B-侧增加 C$\times$。面缺陷（螺线管／分歧）则是沿一个曲面插入，二者的对偶由 Fourier 级数在“额外 Z”与“额外 C$\times$”之间互换给出，与 (4.19)--(4.21) 中 的 Dirac 量子化是同一回事。
\end{quote}

\section{类域论}
\begin{quote}\small
这里的缩写 CFT 指 class field theory（类域论），但在数学物理 文献中 CFT 通常指 conformal field theory（共形场论）。本书第四章讨论的是 拓扑场论与超对称，本章讨论的是类域论；阅读时务必按上下文区分这两个 CFT。
\end{quote}

\subsection*{非分歧类域论}
\begin{quote}\small
(5.3) 原文印作 “C[C$\ell$(f)]”，其中 f 小写， 明显是 F 的笔误。另外这里 C[-] 是“函数空间”的记号（更准确地说， 是对于有限群 C$\ell$(F) 的函数空间），与 (4.72) 中群代数的含义 相同；两边都是有限维的，因为 C$\ell$(F) 有限。
\end{quote}

\begin{quote}\small
这里“类似物”的层次要对齐：拓扑侧是 $\pi$1$\to$ H1=$\pi$1ab，算术侧是 Gal(Fur/F)$\to$Gal(Fur,ab/F)。注意 (5.9) 已经用了交换化， 所以 (5.10)--(5.11) 的分解正是“特征只看见交换商”这一事实的特例； 反过来，这也说明单靠特征无法恢复非交换的 Gal(Fur/F)。
\end{quote}

\begin{quote}\small
Z(1)=$\mu$n 是 Tate 模，其“取向”相当于选取 一个相容的单位根系统；R 上的对应物是 2$\pi$ iZ。这就是为什么 (4.19) 中的 b$\Sigma$=(1)/(2$\pi$ i)$\int$ F 与 (5.80) 的 Lefschetz 对偶里 都会出现 2$\pi$ i：算术侧用 Z(1) 代替 Z，等价于在拓扑侧把常值系数 换成带 twist 的系数。原文说“$\mu$n with respect to the algebraic closure of Z/pZ”，准确说法应是“相对于 SpecFp 的代数闭包的 p-进 Tate 模”。
\end{quote}

\subsection*{函数域}
\begin{quote}\small
(5.24) 中 WCur 的定义是 WCur:=($\pi$1ét)-1(Z)$\subset$$\pi$1ét， 其中 Z$\subset$Z 是自然包含。这个构造的意义在于：把 Gal(Fq/Fq)$\cong$Z 换成它的稠密子群 Z （Frobenius 生成的子群），就得到“几何基本群 + Frobenius”的半直积结构； 这使 Weil 群成为比 $\pi$1ét 更接近拓扑基本群的对象。
\end{quote}

\begin{quote}\small
定理 2 的准确含义是：Pic(C) 与 $\pi$1ét(C)ab 的投射有限完备化同构，即 Pic(C)$\cong$$\pi$1ét(C)ab；由于 Pic 已是 Z 被有限群 Jac 的扩张，取完备化只是把 Z 换成 Z。这就是为什么 (5.27) 右边写 WCur,ab： WCur,ab 正是以 Z 为商（而非 Z）的那个版本。原文的 “isomorphism on profinite completions” 一句容易误读成“两个群本身同构”， 实际上只是完备化后同构。另外，本章的定理编号是 “定理 2”“定理 3”，与章节号 无关（第 1 章的定理 1 未出现在本章），这是原书的全局编号方式，译文照原样保留。
\end{quote}

\begin{quote}\small
(5.33) 中 $\boxtimes$ 是外张量积：$\chi$ 是 Pic 上的 函数，$\pi$1*$\chi$$\boxtimes$$\pi$2*$\chi$ 是 Pic$\times$ Pic 上的函数。这是“乘法特征”的函子性表述，也可写成 $\mu$*$\chi$=$\chi$$\otimes$$\chi$。对一般（非交换）群，特征只对应一维表示， 这正是为什么下面只考虑交换的 Pic。
\end{quote}

\subsection*{函数域 / Loc}
\begin{quote}\small
(5.42) 的链条是 Z/$\ell$n$\leftarrow$Z$\ell$$\subset$Q$\ell$$\subset$Q$\ell$， 方向很重要：Z$\ell$ 是 Z/$\ell$n 的逆向极限，所以表示沿着反向的 箭头“下降”；而 Q$\ell$ 是 Q$\ell$ 的代数闭包，其上的 连续表示理论（而非任意域同态）才是“与 $\ell$ 无关”的那一套。作者说 “Q$\ell$$\simeq$C 作为域”，其非典范同构确实存在，但它不保持 拓扑（Q$\ell$ 上没有与 C 相容的拓扑），所以不能用它把 $\ell$-进表示直接翻译成复表示；“与 $\ell$ 无关”是更精细的同构陈述。
\end{quote}

\subsection*{函数域 / Pic}
\begin{quote}\small
(5.46) 原文写作 “x$\in$ CZ/unit,k(C)”。 这里的 “unit k(C)” 指 k(C) 的单位群 k(C)$\times$，即主除子所模去的 对象：Pic(C)=x$\in$ CZ/k(C)$\times$（通过 f$\mapsto$$\Sigma$x vx(f)x 把 k(C)$\times$ 嵌入形式和）。作者用 “unit k(C)” 是想强调模去的是函数域的单位（非零函数），但写成商环的 样子容易误解为“Z 模掉一个环”。
\end{quote}

\begin{quote}\small
(5.71)--(5.72) 中 GL1(A) 与 GL1(O) 省略了下标 F，应读作 GL1(AF) 与 GL1(OAF)。另外这里的 $\pi$0 是在拓扑群意义下取的： 由于 R+ 连通而其他因子紧，连通分支数正是由这些 R+ 因子贡献的。
\end{quote}

\begin{quote}\small
(5.76) 原文排作 “=$\simeq$ U(1)r+s-1”，等号与 $\simeq$ 重复， 是排版笔误；这里只保留 $\simeq$。另外严格说连通分支本身是 (R+)r+s-1 的紧商，而 R+$\cong$R 可缩，所以它的同伦型才是 U(1)r+s-1；作者说“连通分支是 U(1)r+s-1”应理解为同伦等价， 这也正是下一句“上同调是秩 r+s-1 的外代数”的依据。
\end{quote}

\subsection*{分歧 / 物理侧}
\begin{quote}\small
(5.80) 的 Lefschetz 对偶通常写成 Hk((M,N),Z)$\cong$ Hn-k(M N,Z)（或带局部系）。 这里 n=3、k=2，右边得到 H1；系数 C$\times$ 来自 Loc1=Hom($\pi$1,C$\times$) 的系数群，与左边的 Z 系数并不严格 对偶（Hom(Z,C$\times$)=C$\times$ 才是正确配对）。所以这个等式应理解为 “Pontrjagin 对偶意义下的配对”而非同调系数上的 Poincaré 对偶。
\end{quote}

\subsection*{分歧 / 算术情形}
\begin{quote}\small
(5.86) 中的 $\Gamma$ 是“带 Frobenius 作用的惯性群”的完备化： m 生成惯性（即绕去心点的单值性）的 pro-$\ell$ 部分，F 是 Frobenius 的 提升，关系 FmF-1=mq 表示 Frobenius 在惯性上通过 q 次幂作用。 这个群正是局部域 K 的 Weil--Deligne 群的 tame 商；它的交换化 $\Gamma$ab$\cong$Fq$\times$$\times$Z 出现在 (5.100) 中。 注意这里的 F 同时被用作域 F 与 Frobenius，原文没有区分，容易混淆。
\end{quote}

\begin{quote}\small
(5.94) 原文顶部节点写作 “GL1(unit,F) 1(AC)”，其中 “unit F” 在书中没有定义， 从上下文（“无穷阶平凡化处处”）与 (5.95) 的商结构看，它应指 GL1(F)（或 F$\times$）而非某个 “unit” 对象。这里照原文 保留 “unit F” 的字样，并在批注中指出。图的箭头关系是： 顶部 $\to$ Pic（竖直）、顶部 (C,S)（斜向）、 Pic(C,S)（斜向）；此处用 array 重排， 箭头方向与原图一致。
\end{quote}

\begin{quote}\small
(5.102) 与 (5.103) 的差别在于是否写出 “ab”：局部类域论的 标准陈述是 K$\times$$\cong$Gal(Kab/K)，而 (5.103) 写成 GalKx/Kx（略去 ab）只有在取完备化/取交换商的意义下才 成立。作者的意图是：(5.102) 是带 ab 的精确陈述，(5.103) 是同一事实的 简写；K 表示投射有限完备化。
\end{quote}

\begin{quote}\small
(5.104) 与 (5.89) 是同一个对偶的两次陈述：物理侧是 “挖去 $\gamma$ 后的 Loc” vs “沿 $\gamma$ 平凡化的 Pic”， 算术侧是 “在 S 处温和分歧的局部系统” vs “在 S 上平凡化的线丛”。 关键在于两者的“额外因子”必须匹配：Loc 一侧每去掉一个点就多一个 C$\times$（或 Q$\ell$$\times$），而 Pic 一侧每平凡化 一个点就多一个 Z（或 Fq$\times$$\times$Z），这由 Fourier 变换 在 Z 与 C$\times$ 之间互换。
\end{quote}

\section{扩展 TFT}
\subsection*{分歧}
\begin{quote}\small
这里及式 (6.5) 中 $\Gamma$ 的第一个变元其实是场空间 F(M)，而不是流形 M 本身； 作者对 M 与 F(M) 的记号有混用。严格地说，(6.5) 应写作 $\Gamma$(F(M),$\pi$N*EN)， 其中 $\pi$N:F(M)$\to$F($\partial$ M) 是限制映射。后文 (6.19)--(6.21) 改用 F($\cdot$) 才是准确写法。
\end{quote}

\subsection*{分歧 / 拓扑 Maxwell 理论}
\begin{quote}\small
式 (6.18) 在原文中是不完整的：左上角只印出一个方框占位符，右下角写作 Maps(M,K(Z,2))。按上下文（相对上同调 = 纤维积的 $\pi$0），左上应为纤维积 X=pt$\times$Map(N,K(Z,2))Map(M,K(Z,2))， 右下应为 Map(N,K(Z,2))（限制映射的靶）。这里照原文保留方框， 以便读者核对源文；理解时请按上述纤维积补全。
\end{quote}

\subsection*{扩展 TFT / Atiyah--Segal TFT}
\begin{quote}\small
式 (6.25) 中的 $\langle$-,-$\rangle$Z(N) 是范畴中的 Hom/张量配对（在 2-范畴里即 “沿 Z(N) 的态射空间”），不是 Hilbert 空间的内积。若把 Z(N) 想成 Shv(F(N))， 则这个配对就是层范畴中的 $\otimes$Shv(F(N)) 与 Hom。作者用尖括号是为了提示与 向量空间内积的类比，但切勿按内积（共轭线性）理解。
\end{quote}

\begin{quote}\small
式 (6.27)--(6.29) 是“去环路”条件的标准表述，但 2-范畴 C 的环路 EndC(1C) 只在 C 是单对象 2-范畴（即带 E1 结构的幺半范畴）时才是通常意义的 幺半范畴。此处作者取 n=2 的例子（Bord0,2）作直观；对一般的 n，Bordn-2,n 是 n-范畴，去环路需迭代。严格的完全扩展 TFT 由配边假设（cobordism hypothesis， Lurie）刻画：Z 由 Z(pt) 作为可对偶对象唯一决定。建议对照 [Lur09] 第 2.3 节。
\end{quote}

\subsection*{阿贝尔对偶的例子 / 二维情形}
\begin{quote}\small
(6.36)--(6.37) 是本章的关键机制，也是初学者最容易出错的地方： LocC$\times$(pt) 作为集合只是一个点，但作为栈它是 BC$\times$=pt/C$\times$， 其上的拟凝聚层/局部系统给出 Rep(C$\times$)$\simeq$VectZ-gr。 若在 A-侧同样只取同构类，就丢掉了 Z-分次结构，两边不可能等同。 另外需注意：Shv(Z) 这里指 Z（离散群）上局部系统的范畴，等价于分次向量空间， 而不是 Z 上层的导出范畴。
\end{quote}

\begin{quote}\small
原文第 6.3 节的两个子节标题都写作 “Two-dimensions”，这是作者的重复， 并非排版偶然（源文 p.79 与 p.80 均如此）。这里按原文照译，便于对照； 内容上两小节分别处理 H0$\leftrightarrow$ H1 与 H1$\leftrightarrow$ H0 两组对偶。
\end{quote}

\subsection*{四维 Maxwell 理论 / B-侧}
\begin{quote}\small
式 (6.56) 的长正合列容易看漏一处：H1($\Sigma$,O)/H1($\Sigma$,Z) 中的 H1($\Sigma$,Z) 通过指数正合列嵌入 H1($\Sigma$,O)，这一步用的是 H2($\Sigma$,Z)=Z 与 H2($\Sigma$,O)=0（对紧 Riemann 面）。于是 Jac$\simeq$ K(H1($\Sigma$,Z),1)， 其复结构来自 H1($\Sigma$,O)/H1($\Sigma$,Z)。另注意 (6.66) 中对偶环面 TC 的 下标是复数域：这里的“Fourier 对偶”是在复环面层面，$\pi$1 与 H1 互换。
\end{quote}

\subsection*{四维 Maxwell 理论 / 带分歧的电磁对偶}
\begin{quote}\small
式 (6.74) 的记号与上下文不完全一致：左端的 M03 应即三维流形 M0（上标 3 表示维数），右端的 LocC$\times$M 按限制映射 (6.71) 的源应为 LocC$\times$M0 （去掉纽结后的流形），而不是整个 M。作者混用了 M 与 M0。理解时以 (6.73) 为准： 两边分别是 $\Gamma$(Pic(M0),$\pi$*E) 与 $\Gamma$(M0,$\pi$*E)。
\end{quote}

\subsection*{带分歧的类域论}
\begin{quote}\small
原文 (6.81)--(6.86) 有重复排版：式 (6.81) 与 (6.82) 的左端表达式相同， (6.84) 与 (6.82)、(6.85) 与 (6.83) 分别重复（源文 p.84--85）。这里照原文保留， 只为便于对照；逻辑上只需 (6.81)--(6.83) 与 (6.86)。另外 (6.79) 中的 $\prod$' 表示 限制直积：除有限多个因子外，Ky$\times$ 的因子都落在 Oy$\times$ 中。
\end{quote}

\begin{quote}\small
式 (6.98) 末项在原文中写作 Loc1(C)（源文 p.86）。按上下文，这里应是 Kx 的局部域对应，即 Dx$\times$（或 C x）上秩 1 局部系统的范畴 Loc1(Dx$\times$)=Loc1(C x)，而不是整条曲线 C 上的 Loc1(C)。 这可能是作者把 C 与 C x（或 D$\times$）混写；同类混写在 (6.74) 中也出现过。 另外 (6.95) 的 Galab,ur 是整体类域论中的“非分歧 Abel 化 Galois 群”， 而 (6.96)--(6.97) 是局部对应，二者不可混用。
\end{quote}

\subsection*{数域的分歧}
\begin{quote}\small
原文 (6.103)--(6.104) 中 S 的含意需要小心：S 是“允许分歧”的有限位集， 故在 S 外取整（$\prod$v SO$\times$），而作用来自 S 内的 Kv$\times$。 注意 (6.103) 中 $\prod$v 的上标 '（限制直积）在原文 p.87 被省略，严格写法应为 F$\times$$\prod$'v K$\times$/$\prod$v SO$\times$。 此外 p.87 的原文把 “E.g. think Kv=Qp” 插在句子中间，导致 “representations of K$\times$ with modules for functions on GL1 Galois representations” 一句 语序破碎；这里按数学含义译出，但源文的断句确实可疑。
\end{quote}

\section{几何类域论}
\begin{quote}\small
式 (7.3) 中 H1($\Sigma$,Z)=$\pi$1ab($\Sigma$) 用的是 Hurewicz 定理； 而 $\pi$1(Jac)$\simeq$ H1($\Sigma$,Z) 需要先知道 Jac=V/H1 是环面。注意 (7.2) 里 V/H1($\Sigma$,Z) 的 H1 实为 H1 的格（原文写作上同调）， 且 V 是 H0($\Sigma$,$\Omega$1)；两者通过 Poincaré 对偶等同，但符号上不要混淆 H1（格）与 H1（对偶格）。
\end{quote}

\subsection*{Albanese 性质}
\begin{quote}\small
Remark 30 里“Albanese 的基本群是 H1，故对偶的 Pic0 的 基本群是 H1”这句话容易读反：Alb(X) 与 Pic0(X) 互为对偶 阿贝尔簇，对偶把 H1 与 H1 互换，而不是说二者相等。对曲线， Alb($\Sigma$)=Jac$\Sigma$ 的 $\pi$1 是 H1($\Sigma$,Z)， Pic0 的 $\pi$1 则是对偶格 H1($\Sigma$,Z)。另 (7.11) 中的虚线箭头 表示泛性质给出的映射；原文画作从 Pic 到 A 的虚线。
\end{quote}

\subsection*{Dold--Thom 定理}
\begin{quote}\small
(7.15) 里的 Riemann--Roch 需要补充前提：L 的次数为 d，且 d>2g-2 时 h1(L)=0，于是 h0(L)=d-g+1，纤维 P(H0(L)) 的维数是 d-g，故写作 Pd-g 丛。作者说“这些射影空间单连通，所以 AJ• 在 $\pi$1 上同构”，严格论证需要纤维连通（或单连通）且底空间局部平凡的长正合列；这里单连通性 只用到了 Pn 单连通这一事实。另注意 Dold--Thom 定理的精确陈述是 $\pi$i(Sym$\infty$(X,x))$\cong$ Hi(X;Z)（i$\ge$ 1），且需 (X,x) 为 CW 复形；代数几何类比中“群完备化”对应 Sym• 的 Grothendieck 群。
\end{quote}

\subsection*{回顾：非分歧 CFT}
\begin{quote}\small
(7.18) 与 (7.19) 在原文中把同一个等式写了两遍（源文 p.91），这是重复， 这里照原文保留。另需注意 $\chi$E(O(x)) 是 E 在闭点 x 的 Frobenius 上的取值， 严格说需要先固定 E 在基点的纤维（因为局部系统只在同构意义下确定）， 这就是为什么完整陈述要取 C 上的“点+平凡化”或采用 $\pi$1 的表示。 此外 O(x)=AJ(x) 只在 Pic1 的模型下成立。
\end{quote}

\subsection*{特征层}
\begin{quote}\small
式 (7.31) 原文用右箭头 “$\to$” 把 L$\chi$|x 与 $\chi$(L) 连接起来， 但由 (7.30) 应读作张量积：$\chi$(L(x))$\simeq$ L$\chi$|x$\otimes$$\chi$(L)。 这是原文的笔误（源文 p.92），此处照原文保留箭头，理解时按张量积。 另 (7.26) 中的 Tot($\chi$)$\times$ 是把 k$\times$-torsor 去掉零截面 得到的 k$\times$-丛；说 A$\chi$ 是阿贝尔群即该 torsor 有群结构， 等价于 $\chi$ 给出 Pic 到 BGm 的同态，即 Ext1 类。
\end{quote}

\subsection*{几何 CFT}
\begin{quote}\small
(7.42) 原文写作 “S$\subset$$\sigma$”（源文 p.93），其中 $\sigma$ 应为 C， 这是把曲线记号 C 误写成 $\sigma$（也可能与某个子集记号混用）。此处按数学含义译作 S$\subset$ C。 另证明中 “Pd-g 是 Symdd 的纤维” 依赖 Riemann--Roch（d>2g-2 时 h1=0），且“单连通”指的是纤维 Pd-g 单连通； 严格地说，还需要 Symdd 在 d0 时是纤维丛， 这由 Riemann--Roch 保证。最后，“谱分解”一词在此是类比：几何 CFT 把 Loc(Pic) 分解为沿特征层的直和，正如调和分析把函数空间按特征分解。
\end{quote}

\section{谱分解 2.0：Tannakian 重构}
\subsection*{表示范畴}
\begin{quote}\small
这里“半单 $\Leftrightarrow$ 约化”的判据只在特征 0（更一般地，在约化群的表示完全可约的设定）下成立；正特征时Rep G的半单性对应更强的条件（G的约化性不足以保证完全可约，例如SL2在特征 2 下的表示）。作者此处以特征 0 为前提，读者勿将其当作一般定理。
\end{quote}

\subsection*{Fourier--Mukai 理论 / Betti CFT}
\begin{quote}\small
这里把“次数”与“自同构”互换的做法（(8.75)）是理解本段的关键：Pic上的 Hecke 修正（改变次数，即乘以Z x）在对偶侧变成与Wx作张量（改变Gm-稳定化子的作用），二者互为 Fourier--Mukai 对应。核对符号时注意Wx是到分类栈BGm的映射而非函数。
\end{quote}

\subsection*{Fourier--Mukai 理论 / de Rham CFT}
\begin{quote}\small
(8.96)–(8.98) 的对偶计算用到了“对偶把复形反向并移位”的约定： V$\to$ V的导出对偶是V*$\to$V*，再加上[1]恰好把它移回正向。若采用相反的同调/上同调约定，中间一步的[1]位置会移动；这是核对此处符号时最易出错之处。
\end{quote}

\begin{quote}\small
(8.106) 原文写作“ DA1-mod =L| L$\in$ DA1”，左右两端并不真正相等（左边是整个模范畴），应理解为“M是由生成的循环模，即 DA1”。此外原文“for AA1”亦为笔误，应为“for the Weyl algebra DA1”。此处照译并加注。
\end{quote}

\begin{quote}\small
(8.109) 原文把Jac按形式群取商，但严格地说Jac在单位元处的形式群是(A1)g（g为亏格），原文记号从简。此处照译，读者应按AdR=A/ A理解。
\end{quote}

\subsection*{Betti 与 de Rham 的过渡 / $\C^{\times}$的 Betti 与 de Rham}
\begin{quote}\small
注意此处 Betti 与 de Rham 的对照点：Betti 侧允许无限秩对象（如*C），而 de Rham 侧的 D-模都是“小”的、由有限数据生成；反之 de Rham 侧有支撑在点上的对象（如$\delta$1），Betti 侧看不到。两类范畴只在“有限秩/天穹层”这一公共部分上一致，见 (8.131)。
\end{quote}

\subsection*{几何 CFT}
\begin{quote}\small
“不能写成环的同构”的障碍在于Conn1 C只是仿射丛而非仿射簇，因此两侧都不是仿射的；这正是 Betti 情形（Loc1 C为仿射簇）与 de Rham 情形的本质差别。作者在 (8.146)–(8.147) 之后用纤维维数g再次强调了这一点。
\end{quote}

\begin{quote}\small
注意Conn1 C与Loc1(C)虽然作为实流形都同胚于(S1)2g$\times$纤维，但前者是Cg上的仿射丛（复结构、代数几何对象），后者是(C$\times$)g（仿射簇）。作者借此说明 de Rham 与 Betti 两侧的几何性质相反。
\end{quote}

\begin{quote}\small
(8.152) 的交换关系指标不自洽：左边只有xi,$\partial$xi，右边却出现$\delta$ij。正确写法应为$\partial$xixj-xj$\partial$xi=$\delta$ij$\hbar$（i,j独立取值），这是 Weyl 代数的标准关系；照译并加注。
\end{quote}

\section{非阿贝尔}
\subsection*{自守形式 / 自守函数}
\begin{quote}\small
原文此处误写作Bung(C)(Fq)（下标小写g），应为BunG。另注意“非分歧”指只允许在C的有限多个点上分歧；一般情形还需要在指定分歧除子处引入 level 结构，参见 9.1.3 节末的 level N 版本。
\end{quote}

\subsection*{自守形式 / Weil 描述}
\begin{quote}\small
(9.33)–(9.34) 中上标“$\times$ si”是原文记号，读作“取si个拷贝”（即(R$\times$)s1$\times$(C$\times$)s2等）。这里s1是F的实素位个数，s2是复素位个数，故[F:Q]=s1+2s2。
\end{quote}

\subsection*{自守形式 / 上同调自守形式}
\begin{quote}\small
(9.37)–(9.38) 中的G(Z)应理解为$\Gamma$=G(OF)：对一般数域F，整环是OF，G(Z)只是F=Q时的写法。作者在此沿用Q的记号。
\end{quote}

\subsection*{自守形式 / 几何自守形式}
\begin{quote}\small
原文此处在“This is where we will use ideas from TFT.”之后标有橙色批注“transition”，提示讲次过渡；本译本以讲次标记保留该信息。
\end{quote}

\subsection*{Hecke 代数}
\begin{quote}\small
(9.51) 原文印作“C G$\otimes$C K$\otimes$,C”，多了一个$\otimes$，应为C G$\otimes$C KC；此处已按数学含义排版，特此说明。
\end{quote}

\subsection*{Hecke 代数 / 例子}
\begin{quote}\small
(9.91) 的核对方式：T*T在“相等”点处的值等于从该点出发满足“中间点与两端均不等”的中间点个数，即\#A1=q；在“不等”点处的值则等于中间点与其中一端相等的个数，即\#Gm=q-1。于是T2=(q-1)T+q$\cdot$ 1，与 (9.94) 一致。
\end{quote}

\subsection*{Hecke 代数 / 我们情形中的 Hecke 代数}
\begin{quote}\small
末句是本章的悬念，也是后续各章的主题：局部 Hecke 代数Cc[G(O) G(K)/G(O)]的交换性由 Satake 同构（定理 8）保证，但“整体”Hecke 代数的交换性需要 TFT/几何 Langlands 的解释。阅读时勿把两个层级的交换性混为一谈。
\end{quote}

\section{拓扑场论中的算子与缺陷}
\subsection*{态空间}
\begin{quote}\small
这里（10.11）下方的小弧箭头即 id；Zorro 关系就是“拉直配边”给出的复合等于恒等。方向务必核对：求值 V$\otimes$ V'$\to$C、余求值 C$\to$ V$\otimes$ V'，两者顺序不能交换，否则得到的不是对偶公理而是它的转置。
\end{quote}

\begin{quote}\small
（10.12）的四个箭头分别是：上边 ev、左边 Tr（从 V$\otimes$ V* 到 C）、右边 id（从 C 到 End(V)）、下边 $\Sigma$i ei$\otimes$ ei。注意左右两条竖边的方向不同，容易画反；这里 ei 与 ei 是对偶基，Tr(id)=dim V 正是“对偶配对 + 缩并”给出的维数。
\end{quote}

\subsection*{算子}
\begin{quote}\small
原文此处写“因为 A 把它送到一个乘法映射”，其中 A 是上一节的可观测量代数；按上下文这里应为 Z，即函子把配边送到乘法映射。记号混用是原稿的笔误，照译并在此指出。
\end{quote}

\subsection*{界面}
\begin{quote}\small
界面方向的约定最容易出错：O: W$\to$ V 表示“把蓝（W）接到红（V）”，与 (10.13) 中“正点 $\sqcup$ 负点”的解读一致；若把 V 与 W 的位置对调，得到的就是 OT: V*$\to$ W* 而不是同一个界面。另外 (10.21) 中“平凡理论在左侧”对应 C$\to$ V 的箭头方向，画图时不要把它画成 V$\to$C。
\end{quote}

\subsection*{二维 TFT / 基本设定}
\begin{quote}\small
这里 CatC 是 C-线性范畴的 2-范畴；(10.29) 的三角形表示“Bord2 可以提升为取值在 CatC 中的函子”，提升的构造就是 (10.28)：C$\mapsto$Hom(1,C)。注意 Hom(1,C) 是范畴而不是集合，这是 2-范畴与普通范畴的关键差别。
\end{quote}

\subsection*{二维 TFT / 界面与高阶范畴}
\begin{quote}\small
这里是全章第一个容易混淆的地方：Z(pt) 可能不是范畴，但它的“边界条件”却按定义构成范畴 Hom(1,Z(pt))。“界面之间的界面”之所以存在，是因为余维 1 子流形的链环是 S0（两个点），两个端点可以各自独立地选择边界条件。(10.34) 中的颜色正是标记两个端点所接的边界条件，复合则由 (10.37) 的“马裤对”配边给出。
\end{quote}

\subsection*{二维 TFT / 有限群的二维 Yang--Mills 理论}
\begin{quote}\small
本节有三处约定需要记牢。(1) 边界条件的颜色：红/蓝对应 V/W，界面方向是 HomRep(G)(V,W)，即 (10.53) 的 Z(红$\to$蓝)，不要与反方向的 Hom(W,V) 混淆。(2) 正则表示 C[G] 对应 Dirichlet 边界条件、平凡表示 C 对应 Neumann，这与下一节 X=G/1 对 Dirichlet、X=pt=G/G 对 Neumann 的“几何”说法是同一件事。(3) 柱面上的元素是特征 $\chi$V，而不是维数 dim V；维数只在把柱面闭合（再取迹）时出现。
\end{quote}

\subsection*{拉格朗日理论中的边界条件}
\begin{quote}\small
(10.55) 中原文写作 |$\partial$$\sigma$$\subset$ Z，下标是小写 $\sigma$，而全章其余地方用 $\Sigma$；这应是原稿笔误，照译于此指出。另一个更实质的易错点是“边界条件”的措辞：如作者强调的，它其实是“把映射在边界上提升到 Z”的额外结构，而不是对场施加的等式约束；这正是 X=G/1 给出 Dirichlet（平凡化）、X=pt 给出 Neumann（无附加数据）的原因。
\end{quote}

\subsection*{拉格朗日理论中的边界条件 / 边界条件之间界面的复合}
\begin{quote}\small
Example 55 中原文写“X/G=pt/H”，按上下文应为 X1/G=pt/H（X2/G=pt/K 才是另一个），这是原稿的指标笔误，照译并指出。另需注意 K G/K 的复合正是卷积：K=1 时退化为群代数 C G，K=G 时退化为 C，这两个极端之间才是真正非交换的 Hecke 代数。此外 (10.76) 原图为与 (10.74) 同样的“马裤对”配边，仅把三个边界标号都改成 K。
\end{quote}

\begin{quote}\small
(10.89) 原文写作“=C[G]-mod-”，末尾的短横线表示右侧还有一个作用未写出，完整写法应为 C[G]-mod-C[G]（即 C[G]-双模）；这是原文的省略式写法。真正关键的一步是 C[G$\times$ G]=C[G]$\otimes$C[G] 与 C[G]$\simeq$C[G]op（g$\mapsto$ g-1）给出 (10.87)--(10.89) 的链，从而把“乘积理论”翻译成“界面”。
\end{quote}

\subsection*{拉格朗日理论中的边界条件 / 物理学中的局部算子}
\begin{quote}\small
(10.98) 中的积分是形式记号：TFT 中并无真正的测度 D，作者只是借用物理写法来表达“对去心邻域的场求和”。真正严格的表述是 (10.100)--(10.101)：局部算子就是 Z(Sn-1)，而期望值是从 Sn-1 到 $\emptyset$ 的配边所给的线性泛函。另外 (10.104) 的 $\delta$$\gamma$ 是共轭类 $\gamma$ 上的 delta 函数，在 G 有限时 $\langle$$\delta$$\gamma$$\rangle$ 就是一个加权计数，依赖 $\Sigma$ 的拓扑（亏格）。
\end{quote}

\subsection*{局部算子}
\begin{quote}\small
原文 Figure 10.4 只有图题、没有图形（p151 留白），疑为原稿遗漏；本译本照译图题而绝不伪造图形。该图本应画出“挖去三个 3-球”的情形，用来看乘积的结合性：中间的球面把两个小球隔开，从而两种“先乘哪一对”的顺序给出同一条配边。下文的 (10.108)--(10.110) 正是这一结合性的二维图示版本。
\end{quote}

\subsection*{局部算子 / 代数结构}
\begin{quote}\small
结合性的两张图不要“从字面上”看混：(10.108) 是“先粘左边一对、再粘第三个”，(10.110) 是“先粘右边一对、再粘第一个”，两者都同伦于 (10.109)（一个圆盘挖去三个小圆盘），所以乘积结合。交换性则是另一回事：它需要两点在构型空间中可以互换，这只有当 n$\ge$2（此时 $\pi$0 为单点，甚至 n$\ge$3 时基本群平凡）才成立。
\end{quote}

\subsection*{局部算子 / 谱分解}
\begin{quote}\small
(10.112) 的 Spec 里是交换代数 (Z(Sn-1),*)，而不是某个向量空间；MZ0 是“真空模空间”，它未必光滑、未必约化。(10.116) 的因子化同调 $\int$M 不是普通积分，而是把每一点的局部算子代数沿着 M 装配起来的余极限；当 M 连通时它常退化为 Z(Sn-1)，这正是 (10.115) 中 QC(M0Z)=Z(Sn-1)-mod 的来源。
\end{quote}

\subsection*{局部算子 / 二维 TFT}
\begin{quote}\small
(10.136) 左边原文印作 ZG($\Sigma$G)，上标是大写 G，按意义应为亏格 g（右边写的是 $\Sigma$g 与 2g-2），这是原稿的笔误，照译并指出。另外“Rep(G) 在 G 上成层”并不是典范的：如作者所述，要把它变成真正的 QC(G)，需要选代表元；实表示的情形只能得到 gerbe，这正是物理上“反常”的几何来源。
\end{quote}

\subsection*{B-模型}
\begin{quote}\small
本节要抓住“比较”的逻辑：任给交换环 R 都有截断 B-模型 BnX（局部常值映射），它是“最简单的拉格朗日理论”；对任意 TFT Z，用它的局部算子环造出 Bn+1M0Z，谱分解的准确内容就是 (10.152)--(10.154)：Z 是该 B-模型的边界条件，而不是说 Z 与它相等。R$\otimes$$\pi$0(M) 是“每个连通分支一份 R”的简写，不要理解成真的张量积。
\end{quote}

\subsection*{缺陷 / 带奇点的配边假设}
\begin{quote}\small
(10.163) 原文写的是 M0Z=(Sm-1)，上标 m 与上下文（n 维理论）不符，应为 Sn-1，这是原稿的笔误。(10.166) 与 (10.170) 的区别正是 $\pi$1 的作用：D2 中两点构型空间的基本群是 Z（绕数），所以交换两次只给出辫关系而不保证恒等；而在 $\ge$3 维，构型空间单连通，平方才是恒等，于是得到对称幺半 E$\infty$。
\end{quote}

\begin{quote}\small
(10.171) 原文本应是一个范畴 (Z(S2),*)，而 (10.172) 才是对它的谱；这里把“取谱”的对象与谱本身分开写更清楚。M1Z0Z 是仿射化，一般不是同构：线算子比局部算子知道严格更多的信息（见 Remark 64）。
\end{quote}

\subsection*{缺陷 / 局部算子与线算子}
\begin{quote}\small
(10.183)--(10.186) 的核心等式链是 Z(Sn-1)=EndZ(Sn-2)(1)：局部算子是线算子范畴中单位的自同态。这就解释了为什么“局部算子环 = 线算子范畴的中心”。取 Spec 之后，1（结构层），End(O)=$\Gamma$(O) 是整体截面，于是 M0Z=Spec$\Gamma$(OM1Z) 恰是仿射化。(10.186) 中 GCG G 是双模记号（群的左、右作用）。
\end{quote}

\begin{quote}\small
(10.212) 中的 M2Z 疑为 M1Z 的笔误：这一段的语境全是线算子的谱 M1Z 与 Z(Sn-2)，而 M2Z 直到 Remark 65 才作为曲面算子的对象正式定义、且并未再使用。照译并指出。另注意 Maplc 与 Mapslc（带基本群胚的局部常值映射）的区别正是 (10.199) 与 (10.202) 的差别。
\end{quote}

\subsection*{缺陷 / 回到拓扑 Maxwell 理论}
\begin{quote}\small
(10.216) 左边原文印作 A(S1)，但线算子的链环在四维中是 S2，且本段讨论的是 S2=$\Sigma$$\cong$P1 上的线算子，故这里应为 A(S2)，是原稿的指标笔误。链路上：Pic(P1)=Z 是离散空间，其上的局部系统 (= Loc(Z)) 就是 Z-分次向量空间，这正是 (10.216) 右边写成 VectZ-gr 的原因；而 QC(Loc1$\Sigma$) 与 Loc(Pic,$\Sigma$) 是同一个范畴的两种描述（Tannakian / Fourier 对偶），这是几何 Langlands 中“电—磁”两侧的关键识别。
\end{quote}

\subsection*{缺陷 / 拓扑 Maxwell 理论}
\begin{quote}\small
本节的记号是全书最密的：Fx 是 x 处的局部域、Ox 是其整数环、Kx 是分式域（“圈群”G(Kx) 即 G 在局部域上的点）；$\prod$' 表示阿黛尔意义下的受限乘积。(10.261) 中 BunG 与 BunG(C) 混用是原文写法（前者指整体模栈）。(10.262) 原文分母印作 $\prod$y$\neq$ xG(Ox)，按意义应为 G(Oy)，这是原稿的下标笔误，照译并指出。(10.273) 原文在两个集合之间漏印了关系符号，本译本按意义补为 =。最后要记住：几何 Langlands 的“自守侧”是 Shv(BunG(C))，“谱侧”由 Hecke 范畴 Sph 作用；Hecke 算子就是物理上的线算子，而线算子交换，这正是 Sph 虽是 Hecke 范畴却交换的原因。
\end{quote}

\section{几何 Satake}
\subsection*{总体思想}
\begin{quote}\small
原文记号X$\supset$ K是“K作用在X上”的简写，并非包含关系；平衡积X$\times$KG/K的严格含义是(X$\times$ G/K)/K，其中k$\cdot$(x,gK)=(xk-1,kgK)。这一约定在 (11.14) 的LG+-丛中反复使用，读者若不熟悉伴随丛的构造，可对照 [Bor91] 或任何主丛教材。
\end{quote}

\subsection*{仿射 Grassmann 流形}
\begin{quote}\small
(11.14) 原文写作HeckeX，而其余各处均为Heckex，是记号笔误（x为曲线上一点，X为曲线/概形），译文保留原式并在此指出。(11.22) 中的并集应理解为colimN，故只能说Grn是 ind-真 ind-概形而非真概形。
\end{quote}

\subsection*{仿射 Grassmann 流形 / $\Gr$的同伦型}
\begin{quote}\small
(11.42) 中的A$\infty$是形式无穷维仿射空间，原文 OCR 只留下 “Z$\times$ d”，此处依 p.184 页面图恢复。(11.43) 的TPic(C)=H1(C,O)是形式切空间，其右侧O(C x) Kx/Ox为“在x外正则的亚纯函数模去整函数”的商，切勿与H1的 Zariski 切空间混淆。
\end{quote}

\subsection*{仿射 Grassmann 流形 / $\Gr$的几何结构}
\begin{quote}\small
(11.54) 的最后一个双射$\Lambda$/W$\leftrightarrow$$\Lambda$+需要先固定一个支配序（dominant chamber）；$\Lambda$+指支配余权。G(O)-轨道与LG--轨道在Gr上的偏序是相反的（见第 11.3 节与图 11.2），这是理解 (11.68) 中“集合相同但偏序反转”的关键。
\end{quote}

\subsection*{$\Gr$中的轨道}
\begin{quote}\small
(11.68) 中第二个双射严格地说应写成LG+ LG/LG+$\leftrightarrow$BunG(P1)的同构类集合，而非BunGP1本身；原文以同一记号兼指模空间与其点集，译文从之。此处的$\Lambda$/W与LG--轨道的对应关系依赖LG- LG/LG+的分解，与LG+-轨道的$\Lambda$+标号相差一个支配序取反。
\end{quote}

\subsection*{$\Gr$上的等变层 / 因子化结构}
\begin{quote}\small
(11.74) 的半单性依赖于系数域特征为0（D-模范畴）或Q$\ell$-层、且轨道维数奇偶性恒定；若在正特征或整系数下，Ext1未必消失。(11.88) 的因子化同构必须在xi$\cap$yi=$\emptyset$时陈述；在对角线附近应改为邻近循环/因子化范畴的陈述，这正是原文随后讨论的内容。
\end{quote}

\subsection*{$\sph$的 Tannakian 重构}
\begin{quote}\small
(11.95) 中的箭头方向是Spec(C,*)/Gal(C,F)，即把 Tannakian 叠实现为“纤维函子的自同构群”的分类空间；原文排版把$\simeq$与箭头叠在一起，译文按语义写成$\to$$\simeq$。(11.100) 的Shv($\Lambda$)$\simeq$Rep(T)是 Fourier 变换（每点取纤维再对格作 Fourier），需要$\Lambda$为有限生成自由交换群。
\end{quote}

\subsection*{抛物诱导/抛物限制 / $\Gr_B$}
\begin{quote}\small
Claim 2 只在k-点（甚至只在代数闭域的点）上成立；GrB$\to$GrG在族中既非单射也非满射，因此不能据此断言二者作为 ind-概形同构。习题 3 中N(K)的轨道S$\lambda$与G(O)-轨道方向相反：前者无穷维、有限维截面的余维无穷，这正是 MV 循环的支撑。
\end{quote}

\subsection*{抛物诱导/抛物限制 / 回到$\Gr_G$上的抛物限制}
\begin{quote}\small
(11.125) 中的H是双曲平面；GR/K与GrG的类比并非同构，而是“非紧对称空间 vs. 仿射 Grassmann 流形”的层次类比。这里NR-轨道（极限圆）对应N(K)-轨道（半无限轨道S$\lambda$），而T-轨道对应不动点格GrT$\simeq$$\Lambda$；Mirkovi'c–Vilonen 的纤维函子正是“沿极限圆积分”的代数版本。
\end{quote}

\subsection*{谱分解 / 本征对象}
\begin{quote}\small
(11.142) 中“$\otimes$标记的箭头”表示这是张量函子（$\otimes$-functor）；Ex之所以由G-torsor 给出，是因为 Tannakian 对偶把Shv(pt)上的张量函子等同于pt/G的一个点（即G-torsor）。原文把G的法向分次（由2$\rho$给出）单独列出，是为了让“本征值”与表示的分次相容。
\end{quote}

\subsection*{谱分解 / 变化点}
\begin{quote}\small
11.7.2 开头原文把右端误写作Rep(G*)（星号），其后一律为G；译文保留该处笔误形式以对应原文，读者应读作G。(11.145) 的Maplc指局部常值映射，其叠结构使Spec只能取导出意义；原文在 Betti 情形下不加导出修饰，读者需自行补上。
\end{quote}

\subsection*{朴素的几何 Langlands 猜想 / 朴素的电磁/Montonen–Olive $S$-对偶}
\begin{quote}\small
(11.152) 按原样不成立，正确的修正通常是把QC换成IndCoh（Arinkin–Gaitsgory 的“修正几何 Langlands 猜想”），因为LocG奇异时QC丢失了导出信息；这也是 (11.165) 只能是包含而非等号的原因。(11.159) 中的Atriv,L依赖 Levi L，不同L给出不同的 Eisenstein 层，但都以平凡局部系统为本征值——这正是“本征值不能唯一确定本征对象”的表现。
\end{quote}

\section{范畴化}
\subsection*{基础}
\begin{quote}\small
这里的CGG是G上的类函数空间（即G在C[G]上的不变量，或等变层Shv(G/G)的截面），并非“G的平方”；K0(Rep(G))与类函数空间在系数域上同构，正是特征标理论的起点。若把G换为有限群，K0(Rep(G))就是表示环。
\end{quote}

\subsection*{（去）范畴化与 TFT}
\begin{quote}\small
(12.2)–(12.9) 的配边图（杯、帽、S1）在原文中是手绘小图，本译本以、、重新排版；+与-表示边界点的定向。Mark of Zorro 条件 (12.19) 只写出了一条 zig-zag 恒等式，另一条(id)$\circ$(coev)=id需对V单独验证，二者合起来才给出End(V)$\simeq$ V$\otimes$ V。
\end{quote}

\subsection*{例子}
\begin{quote}\small
H*(A)在非导出情形等于A$\otimes$A$\otimes$ AopA，其导出版本应取 Tor，得到k[x]与外代数$\Lambda$•(k[x])的张量；原文只给出“degree 1 处另一份k[x]”的粗略说法。(12.33) 的同构一般需要K0取在足够好的范畴上且k包含足够多的单位根，有限群情形要小心特征p。
\end{quote}

\subsection*{迹}
\begin{quote}\small
(12.35) 的“$\simeq$”只在导出 Morita 理论中成立；在普通范畴层面，A$\otimes$ Aop-模到End(A-mod)的映射未必满。注意Tr与ev是同一个赋值映射的两种写法：Tr(M)=ev(M)，其中M是M对应的自函子。此外 (12.37) 的dim(A)是 Morita 范畴中的维数（即HH*），与第 12.1 节中用K0(C)定义的维数不是同一个东西（见警告 1）。
\end{quote}

\subsection*{$\Bord$中的$\dim$ / $\Bord$中的迹}
\begin{quote}\small
(12.42) 的两个配边在原文中画成“杯”与“帽”；coev的余定义域应是M$\sqcup$M，不要漏掉反向定向。Bord中“迹=映射环面”与下一节Corr中“迹=不动点”是同一对称幺半函子图景的两端，二者由 (12.63) 的分解联系起来。
\end{quote}

\subsection*{$\dim$的几何/拉格朗日解释 / 与范畴化的关系}
\begin{quote}\small
(12.51) 在朴素集合层面$\Delta$$\cap$$\Delta$=$\Delta$，真正的信息在导出纤维积：对 orbifold/栈，X$\times$X$\times$ XX记录了对角线的导出自交，这正是“惯性栈”与自由环路空间的区别。(12.67) 的对称幺半性与 §12.6.2 的“迹送到迹”是层论版本的核心假设；在凝聚层或反常层中它通常只在导出意义下成立。
\end{quote}

\subsection*{几何 Satake}
\begin{quote}\small
(12.90) 与 (12.93) 中Sph（代数）与（范畴）用不同字体区分，原文用粗体与花体混排、容易混淆；这里统一以Sph记球 Hecke 代数、记球 Hecke 范畴。(12.97) 的L是惯性/自由环路空间；对pt/G它等于G/G，故QC的维数给出G的表示环而非群本身——这正是 Tannakian 重构（保留张量结构）比K0（只看表示环）更强的地方。
\end{quote}

\section{Langlands 对应}
\subsection*{几何 Satake}
\begin{quote}\small
源文此处写作“AG(S2$\times$ S2) 的元素就是圈的可能的标号”，与同一句前半的链环 S2$\times$ S1 矛盾，应为 S2$\times$ S1（即 (13.2) 左端）。译文按上下文取 S2$\times$ S1，并在此标明原文笔误。
\end{quote}

\begin{quote}\small
(13.1) 左端的 * 是线性对偶，不是范畴化意义下的对偶；书中用它强调 AG(S2) 是“线算子所成范畴”在某种意义下的 K0 型不变量。(13.2)--(13.3) 用到了配边假设：S2$\times$ S1 是 S2 上的圆族，取“圆上的迹”给出 dim。切勿把 dimAG(S2) 理解为向量空间维数，它是范畴的迹（即去范畴化后的维数）。
\end{quote}

\begin{quote}\small
物理侧的对应是：'t Hooft 算子的电荷原本是 H2(S2;Z)=Z（交换情形），非交换情形应取在 Rep(G) 中；“表示环元素”只是其 K0 影子。真正的标号在导出/范畴化层面是一个 Rep(G)-模结构，这也解释了为什么后面要出现 |Rep(G)| 与 Rep(G) 两个层次。
\end{quote}

\subsection*{函数域上的非分歧 Langlands 对应}
\begin{quote}\small
(13.18) 的交换三角形在原书中以手绘箭头给出（p. 223），此处以 array 重排：底部的张量积同时映到 AutC（经由 Hecke/卷积算子）与 O(Loc)（经由 Wilson 圈）。问题在于这后一个映射是否“穿过” AutC。
\end{quote}

\begin{quote}\small
(13.15) 中把 O(G/G) 与 |Rep(G)| 等同，依赖 G 是约化群、O(G)G=O(G/G) 为不变函数环；这是 GIT 商而非栈商，见 (13.60) 的讨论。
\end{quote}

\begin{quote}\small
(13.19) 用 $\omega$ 而非 O是本章的技术核心之一：在光滑栈上 $\omega$ 是“体积形式/对偶复形”，其函子性（pushforward）比函数环好；若误写成 AutC$\simeq$O(Loc)，则在分歧点处无法得到正确的上同调支撑（参见 (13.117) 中 O-mod $\neq$QC 的现象）。“$\omega$ 版本”也解释了为什么 (13.118) 中要用 IndCoh 而不是 QC。
\end{quote}

\subsection*{分歧 Langlands 对应 / 自守侧的分歧}
\begin{quote}\small
源文此处写作 $\prod$x$\in$ S，与紧接着的“在右边少模掉一些”以及 (13.22) 中的 $\prod$x SG(Ox)$\times$$\prod$x$\in$ SKx 相矛盾：右端应为对全部 x$\in$ C 的乘积，或至少 $\prod$x S。此处照译原文并标明，读者应按 (13.22) 理解。
\end{quote}

\subsection*{分歧 Langlands 对应 / 谱侧}
\begin{quote}\small
(13.29) 左端用的是不带撇的 （原文如此），而 (13.28) 用的是限制张量积 '；两者在“限制”意义下的区别是源文的记号疏漏，不影响实质。
\end{quote}

\begin{quote}\small
(13.37) 与 (13.18)、(13.30) 同型：三角形说明 “大张量积同时作用在自守侧与谱侧”，而 Clozel 猜想断言的是“作用穿过自守侧”——但仅在代数自守形式的全子范畴上。原文图以手绘箭头给出（p. 227），此处重排。
\end{quote}

\begin{quote}\small
代数自守形式（algebraic automorphic forms）与上同调自守形式的区别不容混淆：前者要求无穷素点处的无穷小特征属于某个格（即“代数”），后者要求出现在 H*([G]F) 中；Clozel 猜想通常陈述为上同调自守形式（或代数自守形式的有理结构）对应 Galois 表示。
\end{quote}

\begin{quote}\small
这一段是全书对“动机”最直白的解释：单个 $\ell$-进表示是实现，动机才是“对所有 $\ell$ 一致的母对象”。技术上 H*M(X) 应理解为一个 Chow 型动机（或至少是 numerical/homological 动机）中的对象，其各种实现（de Rham、Betti、$\ell$-进平展）是纤维函子；把它当成真实存在的 Q-向量空间会出错。
\end{quote}

\subsection*{L-函数}
\begin{quote}\small
(13.41) 右端的 在原文中没有下标，应为 (13.39) 中的 ' p S；这是源文记号疏漏。这里照译。
\end{quote}

\begin{quote}\small
(13.47) 与 (13.45) 的对偶是 Koszul 对偶的“多项式层面”痕迹：• V 的分次迹给出 det(1-tF)，而 Sym• V 的分次迹给出 1/det(1-tF)。若要写下解析延拓后的函数方程，须把 s 归一化为 s$\leftrightarrow$ 1-s 的中心化形式，并把 F p 换成酉化（unitary）的 F p|$\cdot$|1/2；书中未做此归一化。
\end{quote}

\begin{quote}\small
(13.53) 中的 p 是 p 的剩余域大小（有理素数情形即 p）；数域情形应写 N( p)-s。此外 (13.56) 与 (13.53) 的 L-因子只有当 $\rho$ 在 p 处非分歧、且 $\rho$ 的像落在 G 中时才自然对应。
\end{quote}

\subsection*{整体 Langlands}
\begin{quote}\small
(13.60) 中 G/!/G 是 GIT 商，其点正是半单共轭类（等价地特征多项式系数）；而 $\lambda$ p 若来自非半单的 F p，就只记住其半单部分。这正是“只用 G/!/G 而不只用栈 (G)/(G)”会丢失信息的地方，也是需要栈化/局部 Langlands 的动机之一。
\end{quote}

\begin{quote}\small
原文称 Langlands 猜想蕴含 L($\rho$,V,s) 对不可约 $\rho$（dim>1）为整函数；严格说这需要 $\rho$ 不是一维的（一维时即 $\zeta$-函数有极点），而且“整”要理解为在适当归一化、去掉可能的极点后的整性。自守侧的解析延拓（即 L(f,V,s) 的整性）由 Langlands–Shahidi 等方法给出，这是“L(f)=L($\rho$)”最有用的推论。
\end{quote}

\subsection*{局部 Langlands 对应 / 一般形态}
\begin{quote}\small
(13.79)--(13.80) 之所以把配边当作“对象”，是因为 n 个圆是输出边界；若把 $\Sigma$0 反过来读作从 S1 到 $\emptyset$ 的配边，它给出的是 AG(S1)$\to$1 的配对（即迹）。书中只取了前一种读法。另注意这里的 TFT 必须已取定切触结构（定向）；见第 14 章 Remark 101。
\end{quote}

\subsection*{局部 Langlands 对应 / 局部范畴}
\begin{quote}\small
(13.90)--(13.91) 右端的 $\prod$ 在原文中没有指标，应理解为 (13.82) 中的 $\prod$ p SG(O p)。这样 (13.90) 才是“在 S 处不模掉、在 S 外模掉 O”的版本，(13.91) 则是 S 处也模掉 $\prod$ 的版本；两者相差的正是 (13.87) 中的嵌入。
\end{quote}

\subsection*{局部 Langlands 对应 / 谱侧}
\begin{quote}\small
(13.98) 中若用 Weil–Deligne 群，则 Gal(K) 应替换为 WK2（或 Weil–Deligne 参数），否则无法同时记录半单部分与单幂部分；书中此处只是给“朴素猜测”，下一段的局部系统计数问题（非半单 Frobenius）正是需要这个修正的地方。
\end{quote}

\begin{quote}\small
(13.99) 的 LocG(D$\times$) 在代数几何中不是通常意义下的栈（形式穿孔圆盘没有足够多的点）；把它当作栈需要取形式/导出意义上的局部系统，或直接在完美胚（perfectoid）框架下处理，见 Remark 97 与 13.6.4 末段。
\end{quote}

\begin{quote}\small
(13.100) 的图形在原书中为手绘示意（p. 237），此处保留原图。其含义是 S2 是 S1 上的“两个圆盘”的粘合，对应把线算子(S2) 写成面算子(S1) 的自同态；这正是 (13.101) 的内容。
\end{quote}

\begin{quote}\small
定理 13 中的“秩 n Galois 表示”应理解为“n 维 Galois 表示”（rank 指 GLn 标准表示的维数，而非概形的秩）；且 LocGLn 此处指 Gal(K) 到 GLn 的连续半单表示（带适当分歧条件）的栈， 表示同构。
\end{quote}

\begin{quote}\small
(13.102) 的 Shv(pt/Z($\rho$)) 按定义就是 Z($\rho$)-表示；之所以只剩下中心化子，是因为在栈上取点 pt$\to$Loc 的纤维积要模掉稳定化子作用。这里 Z($\rho$) 可取为 $\rho$ 的 Gal-中心化子（若 $\rho$ 不可约），其 $\pi$0 部分给出 L-包，连通部分给出无穷小特征。
\end{quote}

\begin{quote}\small
(13.107) 中仿射 Grassmann 流形 T(K)/T(O) 的同构 $\simeq$$\Lambda$ 需固定一个均匀化子 t；作为格它是 $\Lambda$（余特征格，cocharacter lattice），而非特征格。原文称 $\Lambda$ 为 “character lattice”，此处照译但请注意：T(K)/T(O) 自然地是余特征格。这一字之差会影响后续 $\Lambda$-mod 与 QC(T) 的对应方向。
\end{quote}

\begin{quote}\small
(13.108) 中 T 是环面，QC(T) 即 $\Lambda$-分次向量空间（取 $\Lambda$ 为 T 的特征格时）。若按原文把 $\Lambda$ 当特征格，则 QC(T) 应是 $\Lambda$=$\Lambda$ 的分次，二者在环面情形恰好一致，故不影响结论。
\end{quote}

\begin{quote}\small
(13.111) 的 2$\times$2 图中第二行左端 $\Lambda$-mod$\otimes$Rep(T) 是把“稳定化子修正”张量进来；右端 QC(T/T) 正是水平迹（环路空间），它与 QC(T)（竖直迹）的差别就是 (13.116) 的嵌入。(13.113) 中 ($\cdot$,•) 表示分次环（奇偶/度数），而 O(T) 只看平凡分次，故竖直迹丢失了度数信息。
\end{quote}

\begin{quote}\small
(13.117) 的不等号是本章最要紧的一句：O(G/G)-mod 是“朴素迹”，而 QC(G/G)（即水平迹）才与自守侧匹配。这正是 V. Lafforgue 的 excursion 算子与 AGK 的“谱分解”技术的关键差别，也是 (13.118) 要改用 IndCohNilp 的原因。
\end{quote}

\subsection*{局部 Langlands 对应 / 完整/范畴化局部 Langlands}
\begin{quote}\small
(13.118) 的脚注不可略过：光滑/有限型假设下 QC 与 IndCoh 的差别在 BunG 这类非有限型、非光滑对象上是本质的；正确陈述用 IndCohNilp(LocG)。(13.120) 的 Kottwitz 同晶集合并不能直接给出 BunG 的几何结构，取 G(K)/FG(K) 作为“空间”只是对同构类的粗描述。
\end{quote}

\begin{quote}\small
(13.121) 中 F 是 Frobenius 自同构，h*g=hgFh-1 中的 gF 表示 $\sigma$-线性作用；把它读作“h g $\sigma$(h)-1”更不会出错。BunG=G(K)/FG(K) 只对函数域（或完美胚意义下的 Qp）有意义，见 Fargues–Fontaine 曲线。
\end{quote}

\section{导出代数几何}
\subsection*{拓扑场论与 $\mathbb{E}_n$-代数 / 局部算子上的 $\mathbb{E}_n$-代数结构}
\begin{quote}\small
(14.5)--(14.9) 的关键是“构型空间参数化运算”：n=1 时 Confk(D1) 不连通（有 k! 个分量），给出结合的乘法和置换；n>1 时连通，给出交换乘法。所谓“导出”就是把 Confk(Dn) 换成它的链复形/同伦型，从而把“只在同伦意义下交换”精化为 En 结构。切勿把 (14.7) 读成“每个构型给一个映射”，它在导出层面是一个 Confk 上的局部系统。
\end{quote}

\subsection*{拓扑场论与 $\mathbb{E}_n$-代数 / 范畴表现}
\begin{quote}\small
(14.16) 中的 11$\cdots$1 是“一路复合到底”的单个 n-态射的记号，原书用嵌套下标表示（11$\cdot$$\cdot$）；这里按行文线性化为 11$\cdots$1，含义不变：这个自同态代数带有 n 个相互相容的乘法。
\end{quote}

\begin{quote}\small
表 14.1 在原文中是一张 4$\times$4 的排版表（p. 244），图文件 table14-1.png 的裁剪把表后正文一并包含进来、不宜直接插图，故按本书规范（guide §3.3：能可靠重排则优先重排）用 booktabs 重排；内容与原文逐格一致。
\end{quote}

\subsection*{Hochschild 上同调}
\begin{quote}\small
(14.20) 中把中心与 Hochschild 上同调等同，只在 A 为光滑（或至少 HH* 与 EndAe(A) 一致）时成立；一般情形 EndAe(A) 是导出中心，而 Z(A) 只是 $\pi$0。(14.21) 的 BA 指 A-双模。(14.23) 的 Spec 应取导出意义下的谱，这也是本章标题“导出代数几何”的由来。
\end{quote}

\subsection*{Hochschild 上同调 / Gerstenhaber 代数}
\begin{quote}\small
(14.28)--(14.30) 的符号最易出错：Conf2(Rn)$\simeq$ Sn-1 的定向决定了括号的符号（Koszul 符号）；• TM=Sym,T[-1] 中的 [-1] 是“平移”，即把 TM 放在次数 1，于是 kTM 落在次数 k，这与 HH* 的度数一致。(14.33) 的 Spec 是导出 Spec（取分次交换环的谱），所以得到的是 1-移位辛簇而非普通簇。
\end{quote}

\begin{quote}\small
“移位辛”（shifted symplectic）的移位量等于括号的次数：n-移位辛流形的函数环是 Pn（即括号次数 1-n）。(14.32) 的 T*[1]M 是 1-移位辛的，故其函数是 P2（Gerstenhaber）；这正是 HH* 的情形。若把它误当成普通辛流形（0-移位），括号次数会算错 1——这是本节最常见的指标错误。
\end{quote}

\subsection*{量子化}
\begin{quote}\small
(14.36) 的原图（p. 247）是一个小三角形：Assoc 取括号 [,]，Lie 取泛包络代数 U，Assoc 取交换化；此处的线性写法保留了 U 与 [,] 两支，交换化一支图中省略。
\end{quote}

\begin{quote}\small
(14.34) 中分母的 $\hbar$ 与 $\hbar$ 需一起读：ab-ba 在 A$\hbar$=0 处为 0，故可除以 $\hbar$ 后再模 $\hbar$。(14.35) 的 Lie[1-n] 表示括号次数为 1-n 的李代数；n=2 给次数 -1（Gerstenhaber），n=3 给次数 -2，与 (14.28) 的 Hn-1 一致。
\end{quote}

\begin{quote}\small
(14.43) 中 •$\Omega$M[n] 的 [n] 是关键的移位：HKR 同构给出 HHi(M)$\cong$p-q=i$\Omega$pM，作为分次空间 HH*$\cong$•$\Omega$M[n]（每个 $\Omega$1M 移 n 次）。缩并体积形式把 • TM 与 •$\Omega$M[n] 配对，这正是“选体积形式 = 平凡化典范丛 = 给 Calabi–Yau 结构”的来源。
\end{quote}

\section{TFT 中的状态与可观测量}
\begin{quote}\small
Dunn--Lurie 可加性通常表述为 Emnm+n，其几何含义是“n 维小圆盘可以沿 n 个彼此独立的方向分解”；这里作者写成 En=Ek+E$\ell$ 是简写，不要误读为直和。
\end{quote}

\subsection*{几何量子化与形变量子化 / 形变量子化}
\begin{quote}\small
原文（以及抽取文本）中这个形变参数有时被写成 $\hbar$，但页面图 (15.6)--(15.9) 上明确用的是斜体 h，不是 $\hbar$。本译本按页面图统一用 h；这不是把 h 与 Planck 常数 $\hbar$ 混同，而是笔记作者的记号。
\end{quote}

\begin{quote}\small
原文此处人名留空，页面上只有一个手写的“cite”批注框，显系作者尚未补上引用。按数学史，辛流形形变量子化的早期工作应指 Fedosov（以及更早的 Weyl--Wigner--Moyal 传统）；译者不擅自替正文补入人名，请读者据此对照文献。
\end{quote}

\subsection*{几何量子化与形变量子化 / 几何量子化}
\begin{quote}\small
这里作者用“极化”指实极化（把 TCM 分解为全纯与反全纯两部分）；(15.11) 对应取零截面上的函数，是“实极化”的一种。后文 Example 107 说明不同极化给出的 Hilbert 空间应当同构，这正是几何量子化不依赖极化的核心诉求。
\end{quote}

\begin{quote}\small
这里“支撑理论”与波前集、Wigner 分布之间的关系，是后文（15.51）把 Ext(F,F) 看成 T*X[1] 上代数的动机：|f|2 的提升就是通常的“符号”，而奇异支撑可以看作它的代数几何版本。
\end{quote}

\subsection*{（拓扑）量子场论中形变量子化与几何量子化的类比}
\begin{quote}\small
原文句末有手写批注“Cite”，提示作者原打算补一条文献。PTVV13 指 Pantev--Toën--Vaquié--Vezzosi 的 “Shifted symplectic structures”，是该定义的原始文献；Albert Schwarz 则更早用 odd symplectic（即 -1-辛）结构研究 BV 量子化。
\end{quote}

\begin{quote}\small
原文此处印作 “shifted sign symplectic form”，页边还贴了一张写有 “sign” 的批注条。按定义与上下文，这里显然应为 “shifted symplectic form”（移位辛形式），多出的 “sign” 是笔误或批注残留；译文按正确的“移位辛形式”处理。
\end{quote}

\begin{quote}\small
(15.16) 与 (15.17) 的箭头方向值得留意：M0$\to$ DQ Z(Sn-1) 表示“对 M0 作形变量子化得到 Z(Sn-1)”，而不是把代数映回簇。作者用 而非 $\to$，正是为了提示这不是通常的几何态射。
\end{quote}

\begin{quote}\small
(15.22) 中竖直箭头是“代入”映射；交换性就是说，先在边界球面上取定状态 Z(Dn)，再用 Mx 求值，与直接算配分函数 Z(M) 一致。作者原文只画了右侧竖直箭头，这里补成完整的交换方块以便阅读，数学内容不变。
\end{quote}

\begin{quote}\small
更准确地说，因子化同调依赖于流形的切触结构（framing）；对定向 n-流形，需要同伦地选取切标架，SO(n) 作用正是这一选择的残留。严格处理需用 Bordnfr 或相应的 framed 版本。
\end{quote}

\subsection*{凝聚层的奇异支撑}
\begin{quote}\small
对光滑 X，HH*(X)$\cong$ppTX[-p]，这正是对 TX[-1] 取对称代数；对奇异 X，HKR 需要换成导出版本，Sym 应理解为导出对称代数，额外的同调部分会贡献“偶方向”。
\end{quote}

\begin{quote}\small
原文此处有手写批注“cite Beraldo”，指向 Beraldo 关于 1-移位微分算子与奇异支撑的工作（参见其后定理 17）。另外，“1-移位微分算子”不是标准术语，应理解为“量子化 T*X[1] 所得的 E2-代数”。
\end{quote}

\begin{quote}\small
原文此句写作 “T*[1] is set-theoretically the same as X”，漏掉了底空间 X；从上下文（前一句说 T*X[1]）应为 T*X[1]。
\end{quote}

\subsection*{凝聚层的奇异支撑 / 凝聚层的奇异支撑}
\begin{quote}\small
原文此处说 “automorphisms”，但整段讨论的是自同态（endomorphisms）Ext(Ox,Ox)，且下文正是用 Ext 的无限维性刻画奇点；按上下文 “automorphisms” 应为 “endomorphisms” 的笔误。
\end{quote}

\begin{quote}\small
这里“一个方程”指 2g 个元素之积的换位子等于 1；图中竖直箭头 $\Pi$[,] 表示把 (g1,…,g2g) 送到换位子 $\prod$i[gi,gi+1]。LocG(C) 的 lci 性正是指它由这个截面在单位元处的纤维给出。
\end{quote}

\section{（导出）几何 Langlands 中的缺陷}
\subsection*{通过辛几何研究 TFT / 导出版本}
\begin{quote}\small
“可加关系”指把链环维数降低一维时，算子代数降一阶：E4-代数 Z(S3) 正是 E3-范畴 Z(S2) 中单位对象的自同态代数。注意 (16.6) 与 (16.7) 中 Z(S2) 先是范畴、后又作为对象出现，这是笔记的简写；严格写法是对范畴取单位对象 1Z(S2)。
\end{quote}

\subsection*{通过辛几何研究 TFT / 几何 Langlands 中的移位辛空间}
\begin{quote}\small
这里作者把两种描述混在一起：M1BG=g*[2]/G 按定义是 2-移位辛的；而 T*[3](pt/G) 是 3-移位辛的。两者的关系是 g*[2]/G$\simeq$ T*[3](pt/G)：切空间 g[1] 的余切移 [2] 得 g*[2]，而 pt/G 的余切丛本身已带 3-移位（因为 T(pt/G) 在次数 1，余切在 -1，再加 3）。作者所说“-1+3”即指这个抵消。
\end{quote}

\subsection*{通过辛几何研究 TFT / 它从何而来？}
\begin{quote}\small
这里的 pt$\times$gpt 中，g 应理解为 G 在单位元处的形式邻域（或直接取 g 作为加群）；(16.27) 的等同只在单位元的无穷小邻域上成立，这正是“环路非常小”的含义。与 (16.25) 的差别在于把群换成了它的李代数。
\end{quote}

\begin{quote}\small
原文此处页边有手写批注“Good reference: Goresky--Kottwitz--Macpherson”。这批注指向 GKM 关于等变上同调与环面作用的经典工作；在下文的类比中，S=Sym(V[-2]) 正是 S1 等变上同调 H*S1(pt)=k[u]（|u|=2），见 (16.38)。
\end{quote}

\subsection*{局部算子}
\begin{quote}\small
这里用的是 Chevalley 限制定理：(Sym,g)G$\cong$(Sym,h)W，即不变多项式由其在 Cartan 上的限制决定；特征多项式映射 g**/W 正是这一限制的几何实现。
\end{quote}

\subsection*{局部算子 / $B$-侧}
\begin{quote}\small
原文此处有手写批注“cite”，提示补引 Bezrukavnikov--Finkelberg 关于导出几何 Satake 的工作。该定理通常表述为 Shv(LG+ LG/LG+)$\simeq$QC!(LocG(S2))，等价于 g*[2]/G 上的导出凝聚层；右端的 (Sym(g[-2])-mod)G 是对偶描述。
\end{quote}

\subsection*{边界条件（余维一缺陷） / Langlands 函子性}
\begin{quote}\small
原文明明讨论的是 automorphic forms（自守形式），但抽取文本两处都写作 “automorphism forms”（自同构形式），这是笔误；译文按“自守形式”处理。
\end{quote}

\subsection*{边界条件（余维一缺陷） / 哈密顿作用}
\begin{quote}\small
原文此处页边有手写批注“cite BFN and GW”，提示补引 Braverman--Finkelberg--Nakajima（Coulomb 分支）与 Gaiotto--Witten（边界条件与规范理论）的文献。译文正文保留了对这两项工作的提及。
\end{quote}

\begin{quote}\small
原图是一幅不太规整的草图：Sym,g 有斜箭头分别指向 O(T*M) 与 Vect(T*M)，g 有箭头指向 Vect(T*M)。这里按数学含义重排为交换方块：上横为 (16.85) 的 H，右竖为“取哈密顿向量场”（Poisson 括号），下横为 G-作用，左竖为 g,g。
\end{quote}

\begin{quote}\small
原文 “the BG” 指分类栈 BG；“killing form” 应为首字母大写的 Killing 型。此处作者强调：Chern--Simons 的线算子范畴（模张量范畴）是半单的、没有导出结构，因此不能用局部算子/线算子的谱重构出 BG 的 2-移位辛几何；这说明本节的重构程序并非普遍适用。
\end{quote}

\section{几何 Langlands 中的边界条件}
\subsection*{边界条件概述}
\begin{quote}\small
式 (17.1) 中的 H*(G) 作用是 Cartan 模型给出的：G 在 de Rham 复形上的 作用可同伦平凡化，从而在 H*(X) 上只剩 H*(G)（即 G 的链复形/群同调）的作用； 这就是为什么作者称其为“局部常值表示”而不是普通表示。对连通紧群，H*(G)$\cong$ H*($\Omega$ G)， 它与 H*(BG) 通过 Koszul 对偶相关联，于是 (17.2) 中 H*(G)$\otimes$ H*(BG) 正是 G-等变上同调在自由作用情形下的局部模型。注意不要把 H*(G) 与群环 C[G] 混淆。
\end{quote}

\subsection*{相对谱}
\begin{quote}\small
“相对谱”一词的关键在于 (17.9) 中的映射是到 g*/G 而非到点：通常的 Spec 只给出 O(M)G，丢掉等变结构； 相对版本把 T(S2) 看作 QC( g*/G) 中的 E3-代数， 其相对谱保留了矩映射与 G-作用。另外 (17.8) 中省略的“移位 2”是 S2 作为二维流形带来的 cohomological shift；在线算子范畴中它对应 QC( g*/G) 的 2-移位辛结构。
\end{quote}

\subsection*{$L$-函数（B-侧）}
\begin{quote}\small
(17.14)--(17.16) 中“det 的 (-1)k+1 次幂”容易出错：对 k 为偶数 （偶上同调），指数为 -1，故取 det 的倒数；对 k 为奇数，指数为 +1。 这正是超迹 Trgr=Treven-Trodd 在对称/外代数上的体现。另 (17.22) 的 $\zeta$(0)=$\Sigma$ 1=\#N=-1/2 是 $\zeta$ 的 解析延拓值（正则化），并非通常意义下发散级数的和；作者用它是为了强调 L-值需要重正化。
\end{quote}

\subsection*{$L$-函数（B-侧） / 局部故事}
\begin{quote}\small
(17.58) 中 Tr(g,O(M)) 一般没有意义（O(M) 无穷维，且 g 的特征值连续），这正是作者说“有很多极点”的意思：它只是形式上的迹， 真正的对象是 (17.60) 中的 $\mu$*OM。另外 (17.60) 的等号应理解为 G-等变层的（导出）函子化版本，而不是逐点相等；右端的 QC( g*)G 是 G-等变拟凝聚层，对应 g*/G 上的 QC 层。
\end{quote}

\subsection*{$L$-函数（B-侧） / 小结}
\begin{quote}\small
(17.81) 中的 RWM(C) 是 Rozansky--Witten 理论在曲面 C 上 的状态空间（几何量子化），O(Maplc(C,M)) 是因子化同调给出的可观测量。 作者说前者“像旋量模”、后者“像 Clifford 代数”，是因为当 M=T*X 时 O(T*X) 上的辛结构使可观测量呈 Weyl/Clifford 型代数，而模对应极化选择。 另注意 (17.77)--(17.80) 中 Tr(F) 记号省略了所作用的层，应理解为 Frobenius 在 LM 的纤维上的迹。
\end{quote}

\subsection*{关于 $L$-函数的回顾}
\begin{quote}\small
Remark 122 的 “P3-variety” 原文如此（p.291），含义并不标准：按上下文， 应是把 M 的 3-移位 Poisson 结构（P3）量子化为 E3-代数； “取上同调”指忘记 E3 结构只保留交换代数。另 (17.98) 中的 $\omega$ 是 T*X 上的全纯体积形式/Canonical bundle 的平方根型对象，之所以用 IndCoh 是因为 T*X 上自然出现的是（移位）相干层而不是拟凝聚层。
\end{quote}

\subsection*{自守形式的周期（A-侧）}
\begin{quote}\small
(17.109) 末尾的 “/,” 在原文（p.294）中分母留空，应为与 (17.108) 中 “level” 相匹配的水平结构（对 H 取同一个 level 的子群）。作者显然省略了它。 另 (17.101) 的 horocycle 是 N(Z) N(R)（即 y=1 的轨道），其积分给出 q-展开的常数项 a0；而 a1 是 Whittaker 周期，对应 N-周期取非平凡特征。 (17.106) 的推导需要 是模形式，且 an 按 q=e2$\pi$ i z 展开。
\end{quote}

\subsection*{自守形式的周期（A-侧） / Iwasawa--Tate}
\begin{quote}\small
(17.116)--(17.120) 的关键是把 Z 上的几何级数 $\Sigma$$\chi$vn 写成 Kv$\times$ 上的积分：n 对应 n$\in$Z，而 1Ov 是“基本函数”。 这正是 Iwasawa--Tate 理论的核心：L-因子是局部基本函数的 Mellin 变换。 (17.128) 中求和从 $\gamma$$\in$ F 换成 $\gamma$$\in$OF，原因是 $\Phi$$\infty$ 为 Gaussian 而其余 $\Phi$v=1Ov 只在整点上非零（对 F=Q 时 OF=Z）；这一步是 Poisson 求和/函数方程的技术核心，作者未展开。
\end{quote}

\subsection*{一般版本}
\begin{quote}\small
Remark 124 的末两句在原文（p.297）中指代混乱：“The name is motivated by hypertoric varieties. This name comes from them being hyperkähler.”——两个 “name” 指的 应是不同的东西（“multiplicity 1” 的命名动机与 “hypertoric” 的词源）。按数学含义， hypertoric 簇是“超 Kähler 的环面簇”，其名称来自 hyperkähler + toric； 而“重数 1”指的是 O(M)G 在 Poisson 交换意义下的重数。这里照原文译出并提示。
\end{quote}

\begin{quote}\small
(17.148) 左端的 X(K)/X(O) 在原文（p.298）中如此印出，但按上下文 （球函数是 X(K) 上 G(O)-不变的函数）这里应为 X(K)/G(O)，即 “G(O) 作用在 X(K) 上”。分母写成 X(O) 是原文笔误（也可能是把 商 X(K)/G(O) 与 X(O) 的记号混用）。这里照原文保留并在 (17.148) 中提示； (17.149) 起即改用正确的 X(K)/G(O)。
\end{quote}

\begin{quote}\small
(17.148) 与 (17.152) 中的 |Rep(G)| 记号并不标准：作者用它表示 “Hecke 代数所对应的范畴/代数”，即 Satake 等价 k[G(O) G(K)/G(O)
\end{quote}
